\documentclass[11pt]{article}
\usepackage[margin=1.1in]{geometry}
\usepackage{amsmath,amssymb,amsthm,mathtools,booktabs,array,tabularx,longtable}
\usepackage[dvipsnames]{xcolor}
\usepackage[colorlinks=true,linkcolor=MidnightBlue,citecolor=MidnightBlue,urlcolor=MidnightBlue]{hyperref}
\usepackage{fancyhdr}
\pagestyle{fancy}\fancyhf{}\renewcommand{\headrulewidth}{0pt}
\fancyfoot[R]{\thepage}
\fancyfoot[L]{\scriptsize\color{gray}Generated by the Transport Lab, not the version of record. Official version: \href{https://doi.org/10.5281/zenodo.23232104}{doi:10.5281/zenodo.23232104}}
\fancypagestyle{plain}{\fancyhf{}\fancyfoot[R]{\thepage}\fancyfoot[L]{\scriptsize\color{gray}Generated by the Transport Lab, not the version of record. Official version: \href{https://doi.org/10.5281/zenodo.23232104}{doi:10.5281/zenodo.23232104}}\renewcommand{\headrulewidth}{0pt}}
\newtheorem{theorem}{Theorem}[section]
\newtheorem{proposition}[theorem]{Proposition}
\newtheorem{lemma}[theorem]{Lemma}
\newtheorem{corollary}[theorem]{Corollary}
\newtheorem{conjecture}[theorem]{Conjecture}
\newtheorem{refutation}[theorem]{Refutation}
\newtheorem{prediction}[theorem]{Prediction}
\newtheorem{principle}[theorem]{Principle}
\newtheorem{observation}[theorem]{Observation}
\newtheorem{axiom}[theorem]{Axiom}
\theoremstyle{definition}
\newtheorem{definition}[theorem]{Definition}
\newtheorem{fact}[theorem]{Fact}
\newtheorem{assembled}[theorem]{Assembly}
\newtheorem{postulate}[theorem]{Postulate}
\newtheorem{premise}[theorem]{Premise}
\newtheorem{hypothesis}[theorem]{Hypothesis}
\newtheorem{remark}[theorem]{Remark}
\newtheorem{example}[theorem]{Example}
\theoremstyle{plain}
\setlength{\marginparwidth}{0.8in}
\newcommand{\changedmark}[2]{\marginpar{\raggedright\scriptsize\color{#1}#2}}
\providecommand{\Aa}{}\renewcommand{\Aa}{\mathcal{A}}
\providecommand{\Ad}{}\renewcommand{\Ad}{\widetilde{\mathrm{Ad}}}
\providecommand{\Aev}{}\renewcommand{\Aev}{\mathcal{A}^{+}}
\providecommand{\Aod}{}\renewcommand{\Aod}{\mathcal{A}^{-}}
\providecommand{\B}{}\renewcommand{\B}{\mathrm{B}}
\providecommand{\C}{}\renewcommand{\C}{\mathbb{C}}
\providecommand{\CAR}{}\renewcommand{\CAR}{\mathrm{CAR}}
\providecommand{\CS}{}\renewcommand{\CS}{{\mathbb{S}^{+}}}
\providecommand{\CSm}{}\renewcommand{\CSm}{{\mathbb{S}^{-}}}
\providecommand{\Cl}{}\renewcommand{\Cl}{\mathrm{Cl}}
\providecommand{\Con}{}\renewcommand{\Con}{\mathcal{C}}
\providecommand{\Dl}{}\renewcommand{\Dl}{\Delta}
\providecommand{\EN}{}\renewcommand{\EN}{E_{N}}
\providecommand{\End}{}\renewcommand{\End}{\mathrm{End}}
\providecommand{\F}{}\renewcommand{\F}{\mathbb{F}}
\providecommand{\G}{}\renewcommand{\G}{\Gamma}
\providecommand{\HH}{}\renewcommand{\HH}{\mathbb{H}}
\providecommand{\IS}{}\renewcommand{\IS}{I\sigma_3}
\providecommand{\J}{}\renewcommand{\J}{\mathbf{J}}
\providecommand{\LH}{}\renewcommand{\LH}{L_{H}}
\providecommand{\Mat}{}\renewcommand{\Mat}{\mathrm{M}}
\providecommand{\Mo}{}\renewcommand{\Mo}{M^{\circ}}
\providecommand{\Mp}{}\renewcommand{\Mp}{M^{+}}
\providecommand{\Msun}{}\renewcommand{\Msun}{M_{\odot}}
\providecommand{\Mt}{}\renewcommand{\Mt}{\widetilde{M}}
\providecommand{\Ncone}{}\renewcommand{\Ncone}{\mathcal{N}}
\providecommand{\Neff}{}\renewcommand{\Neff}{N_{\mathrm{eff}}}
\providecommand{\PCS}{}\renewcommand{\PCS}{\mathbb{P}{\mathbb{S}^{+}}}
\providecommand{\PT}{}\renewcommand{\PT}{\mathbb{PT}}
\providecommand{\Pin}{}\renewcommand{\Pin}{\mathrm{Pin}}
\providecommand{\Q}{}\renewcommand{\Q}{\mathbb{Q}}
\providecommand{\R}{}\renewcommand{\R}{\mathbb{R}}
\providecommand{\Ric}{}\renewcommand{\Ric}{\mathrm{Ric}}
\providecommand{\SBH}{}\renewcommand{\SBH}{S_{\mathrm{BH}}}
\providecommand{\Sc}{}\renewcommand{\Sc}{\mathbb{S}}
\providecommand{\Sch}{}\renewcommand{\Sch}{\mathrm{P}}
\providecommand{\Sig}{}\renewcommand{\Sig}{\Sigma}
\providecommand{\Som}{}\renewcommand{\Som}{S_\omega}
\providecommand{\Sp}{}\renewcommand{\Sp}{\mathbb{S}}
\providecommand{\Sph}{}\renewcommand{\Sph}{\mathrm{S}}
\providecommand{\Spi}{}\renewcommand{\Spi}{S_\pi}
\providecommand{\Spin}{}\renewcommand{\Spin}{\mathrm{Spin}}
\providecommand{\Sth}{}\renewcommand{\Sth}{S_{\mathrm{th}}}
\providecommand{\T}{}\renewcommand{\T}{\mathbb{T}}
\providecommand{\TH}{}\renewcommand{\TH}{T_{\mathrm{H}}}
\providecommand{\Tr}{}\renewcommand{\Tr}{\operatorname{Tr}}
\providecommand{\UHF}{}\renewcommand{\UHF}{\mathrm{UHF}}
\providecommand{\Z}{}\renewcommand{\Z}{\mathbb{Z}}
\providecommand{\Zc}{}\renewcommand{\Zc}{\mathrm{Z}}
\providecommand{\cmg}{}\renewcommand{\cmg}{\mathrm{cm^{2}\,g^{-1}}}
\providecommand{\dS}{}\renewcommand{\dS}{\mathrm{dS}}
\providecommand{\dsc}{}\renewcommand{\dsc}{\mathrm{dsc}}
\providecommand{\dtPG}{}\renewcommand{\dtPG}{dt}
\providecommand{\dtS}{}\renewcommand{\dtS}{dt_{\mathrm{S}}}
\providecommand{\ellv}{}\renewcommand{\ellv}{\ell}
\providecommand{\er}{}\renewcommand{\er}{e_r}
\providecommand{\et}{}\renewcommand{\et}{e_t}
\providecommand{\ev}{}\renewcommand{\ev}{\mathrm{ev}}
\providecommand{\gl}{}\renewcommand{\gl}{\mathfrak{gl}}
\providecommand{\hatV}{}\renewcommand{\hatV}{\widehat{V}}
\providecommand{\hb}{}\renewcommand{\hb}{\bar{h}}
\providecommand{\hf}{}\renewcommand{\hf}{\hat{f}}
\providecommand{\id}{}\renewcommand{\id}{\mathrm{id}}
\providecommand{\kpc}{}\renewcommand{\kpc}{\mathrm{kpc}}
\providecommand{\lP}{}\renewcommand{\lP}{\ell_{P}}
\providecommand{\mD}{}\renewcommand{\mD}{\mathbf{m}_{D}}
\providecommand{\mb}{}\renewcommand{\mb}{\bar m}
\providecommand{\mnu}{}\renewcommand{\mnu}{\Sigma m_{\nu}}
\providecommand{\nInf}{}\renewcommand{\nInf}{n_\infty}
\providecommand{\nO}{}\renewcommand{\nO}{n_0}
\providecommand{\ninf}{}\renewcommand{\ninf}{n_{\infty}}
\providecommand{\nz}{}\renewcommand{\nz}{n_{0}}
\providecommand{\proj}{}\renewcommand{\proj}{\operatorname{proj}}
\providecommand{\pv}{}\renewcommand{\pv}{\pi}
\providecommand{\rad}{}\renewcommand{\rad}{\mathrm{rad}}
\providecommand{\sgn}{}\renewcommand{\sgn}{\operatorname{sgn}}
\providecommand{\shA}{}\renewcommand{\shA}{\mathcal{A}}
\providecommand{\shC}{}\renewcommand{\shC}{\mathcal{C}}
\providecommand{\shE}{}\renewcommand{\shE}{\mathcal{E}}
\providecommand{\shJ}{}\renewcommand{\shJ}{\mathcal{J}}
\providecommand{\shK}{}\renewcommand{\shK}{\mathcal{K}}
\providecommand{\shL}{}\renewcommand{\shL}{L}
\providecommand{\shW}{}\renewcommand{\shW}{\mathcal{W}}
\providecommand{\sll}{}\renewcommand{\sll}{\mathfrak{sl}}
\providecommand{\so}{}\renewcommand{\so}{\mathfrak{so}}
\providecommand{\su}{}\renewcommand{\su}{\mathfrak{su}}
\providecommand{\supp}{}\renewcommand{\supp}{\operatorname{supp}}
\providecommand{\tr}{}\renewcommand{\tr}{\operatorname{tr}}
\providecommand{\uu}{}\renewcommand{\uu}{\mathfrak{u}}
\providecommand{\br}{}\renewcommand{\br}[1]{\langle #1\rangle}
\providecommand{\bra}{}\renewcommand{\bra}[1]{\langle #1|}
\providecommand{\file}{}\renewcommand{\file}[1]{\texttt{#1}}
\providecommand{\hp}{}\renewcommand{\hp}[2]{( #1 \,|\, #2 )}
\providecommand{\ip}{}\renewcommand{\ip}[2]{\langle #1, #2\rangle}
\providecommand{\ket}{}\renewcommand{\ket}[1]{|#1\rangle}
\providecommand{\rel}{}\renewcommand{\rel}[2]{D\bigl(#1 \,\|\, #2\bigr)}
\providecommand{\rev}{}\renewcommand{\rev}[1]{\widetilde{#1}}
\providecommand{\slashed}{}\renewcommand{\slashed}[1]{{#1}\mkern-10mu/}
\title{\bfseries One Dimension at a Time\\[4pt]\large\itshape The Odd Extensions of the Oriented Towers}
\author{Leon Butler\\ \normalsize\href{mailto:lbutler.math@gmail.com}{\texttt{lbutler.math@gmail.com}}\\ \normalsize ORCID: \href{https://orcid.org/0009-0006-4449-5617}{0009-0006-4449-5617}}
\date{Generated by the Transport Lab, 27 October 2026 \textperiodcentered{} Transport Papers, T14a\\[3pt]\small Official version: \href{https://doi.org/10.5281/zenodo.23232104}{doi:10.5281/zenodo.23232104}}
\begin{document}
\maketitle
\begin{abstract}
\normalsize
The transport adds a null pair and climbs a tower without changing its residue. This note asks what a single added generator does to the oriented towers $\Cl(1,3)$ and $\Cl(3,1)$. A space-sign generator doubles the shadow, $\Cl(1,4) \cong \Cl(1,3) \oplus \Cl(1,3)$; a time-sign generator complexifies it, $\Cl(2,3) \cong \Cl(1,3) \otimes \C$, and there the mirror's algebra sits, on the anti-de Sitter translation bivectors. The first rung's even subalgebra is $\Cl(2,3)$, so the dark-matter paper's copy theorem runs inside the complexification that forgets the real form; the odd extension is the first construction above the shadow, by dimension, that keeps the towers apart, through an element no derived quantity uses. The dark-energy paper's frozen point realises the extension: its scale tractor and Stelle-West vector are the ascent's pair, signed by $\Lambda$, and the algebra of $I^{\perp}$ is the intermediate of this world's own ascent, $\Cl(1,4)$ for the observed $\Lambda > 0$, $\Cl(2,3)$ for $\Lambda < 0$, degenerate at $\Lambda = 0$. By the soldered residue that algebra is this tower's, so $\Lambda$ does not supply the coupling the dark towers need. Run on the degenerate case, the transport returns the rung and a canonical null frame and constrains the frozen point only through its norm: within the series $\Lambda$ is undetermined, its frame components are gauge, and its sign does not select the visible tower. The family over the $\Lambda$ line is a tame wall, completed on both sides in the case the sheaf paper's conservation theorem excludes. The note is a dictionary between three papers, with one new result and no prediction.
\end{abstract}
\medskip
\noindent\textbf{MSC 2020.} Primary 15A66, 83F05; Secondary 83C05, 81R25.\\
\textbf{Keywords.} Clifford algebras in odd dimension; de Sitter and anti-de Sitter algebras; mirror matter; MacDowell-Mansouri gravity; Clifford clock.
\newpage
\noindent \textbf{Authorship and disclosure. }The founding idea of the series is the author's: that division by zero in a Clifford algebra should be read not as a prohibition but as a transport into a higher-dimensional algebra, and that the transported division need not be accompanied by an inverse operation. The direction of investigation and the critical review of the results are likewise the author's. The mathematical development and the drafting of the text were carried out in extended collaboration with Claude Fable 5 (Anthropic), an AI system.

\noindent\fbox{\parbox{\dimexpr\linewidth-2\fboxsep-2\fboxrule}{\textbf{This is not the version of record.} The official version of this paper is the one deposited on Zenodo, \href{https://doi.org/10.5281/zenodo.23232104}{doi:10.5281/zenodo.23232104}. This PDF was generated by the Transport Lab on 27 October 2026 and assembled from approved parts of its store: the results, proofs and text of the paper as the store holds them, each result with its premises and machine checks. It is regenerated when a result it contains is corrected, so it may differ from the deposited paper. Since the deposited version 1.0: 2 results corrected, 3 changed, 17 reworded (wording, notation or citations only) and 37 cited results added. Each revised or added result is marked in the margin (\textcolor{BrickRed}{corrected}, \textcolor{BurntOrange}{changed}, \textcolor{gray}{reworded}), and the last section, \emph{Changes since the deposited version}, lists them with the reason for each.}}
\newpage

\section{Setting, and what is imported}

\phantomsection\label{c:note.odd-extensions.setting-setting-what-imported}
Throughout, $\Cl(p,q)$ is the real Clifford algebra on $p$ generators squaring to $+1$ and $q$ squaring to $-1$, $n = p + q$, with orthonormal generators $e_{1}, \ldots, e_{n}$ and volume element $\omega = e_{1}e_{2}\cdots e_{n}$. The series' convention is time $= $ plus: this world is $\Cl(1,3)$, residue $p - q \equiv 6 \pmod 8$, and its mirror is $\Cl(3,1)$, residue $2$. Under the opposite convention every attribution below to a named member of a reversal pair passes to the other member, and the pair-level statements are unchanged; the dark-matter paper carries the same clause \cite[Thm.~3.4]{ButlerT14}. Nothing in this note depends on which member is called this world except the names.

\begin{assembled}[Imported: the towers and the clock]\label{c:odd-extensions.imported-towers-clock}
The $(1,1)$ ascent $\Cl(p,q) \to \Cl(p+1,q+1) \cong \Mat(2, \Cl(p,q))$ preserves $p - q$; the real Clifford algebras are periodic modulo eight; the eight residues label eight parallel towers. \emph{Provenance:} \cite[Thms.~3.2 and 3.6]{ButlerT3}; \cite[Fact~2.5]{ButlerT14}.
\end{assembled}
\begin{proof}
The $(1,1)$ ascent $\Cl(p,q) \to \Cl(p+1,q+1) \cong \Mat(2, \Cl(p,q))$ preserves $p - q$; the real Clifford algebras are periodic modulo eight; the eight residues label eight parallel towers. \emph{Provenance:} \cite[Thms.~3.2 and 3.6]{ButlerT3} (\ref{c:tower.rigidity-clock-frozen}, \ref{c:tower.commuting-copies}); \cite[Fact~2.5]{ButlerT14} (\ref{c:dark-matter.imported-eight-towers-mod-8-clock}).

\end{proof}

\begin{assembled}[Imported: the representatives at the dimension nearest four, and the types]\label{c:odd-extensions.imported-representatives-dimension-nearest-four-types}
\changedmark{BrickRed}{corrected since the deposit}
At the dimension nearest four the residues are represented by the shadows $(1,3)$ at $r = 6$, $(3,1)$ at $r = 2$, $(2,2)$ at $r = 0$, $(0,4)$ and $(4,0)$ at $r = 4$, $(2,1)$ at $r = 1$, $(1,2)$ at $r = 7$, $(4,1)$ at $r = 3$ and $(1,4)$ at $r = 5$; the residues $0$ and $4$ are sterile and the other six admit matter. The real spinor type is a function of the residue alone: $\Cl(1,3) \cong \Mat(2,\HH)$ (quaternionic, one module) and $\Cl(3,1) \cong \Mat(4,\R)$ (real, one module); in particular the two are not isomorphic as real algebras. \emph{Provenance:} \cite[Prop.~3.1, Thm.~3.4]{ButlerT14}.
\end{assembled}
\begin{proof}
At the dimension nearest four the residues are represented by the shadows $(1,3)$ at $r = 6$, $(3,1)$ at $r = 2$, $(2,2)$ at $r = 0$, $(0,4)$ and $(4,0)$ at $r = 4$, $(2,1)$ at $r = 1$, $(1,2)$ at $r = 7$, $(4,1)$ at $r = 3$ and $(1,4)$ at $r = 5$; the residues $0$ and $4$ are sterile and the other six admit matter. The real spinor type is a function of the residue alone: $\Cl(1,3) \cong \Mat(2,\HH)$ (quaternionic, one module) and $\Cl(3,1) \cong \Mat(4,\R)$ (real, one module); in particular the two are not isomorphic as real algebras. \emph{Provenance:} \cite[Prop.~3.1, Thm.~3.4]{ButlerT14} (\ref{c:dark-matter.towers-are-residues-single-clock-filter}, \ref{c:dark-matter.type-residue-what-each-tower-can}).

\end{proof}

\begin{assembled}[Imported: the copy theorem, and the soldered residue]\label{c:odd-extensions.imported-copy-theorem-soldered-residue}
\changedmark{gray}{reworded}
The mirror tower is a copy of this one: every input to the Standard-Model chain (the conformal spin group $\mathrm{Spin}(2,4) \cong \mathrm{Spin}(4,2)$, the conformal spinor module $\CS$ (twistor module) with its helicity form, the reversion form $h$, the complexified isotropic structure, the mass formula of two conformal spinors) is invariant under signature reversal, because over $\C$ the ambient carries no signature and $\Cl^{0}(2,4) \cong \Cl^{0}(4,2) \cong \Mat(4,\C)$. The residue is a soldered invariant, not an algebraic one, and the one metric superselects the tower charge. \emph{Provenance:} \cite[Thms.~3.5 and 4.3, Prop.~4.1]{ButlerT14}.
\end{assembled}
\begin{proof}
The mirror tower is a copy of this one: every input to the Standard-Model chain (the conformal spin group $\mathrm{Spin}(2,4) \cong \mathrm{Spin}(4,2)$, the conformal spinor module $\CS$ (twistor module) with its helicity form, the reversion form $h$, the complexified isotropic structure, the mass formula of two conformal spinors) is invariant under signature reversal, because over $\C$ the ambient carries no signature and $\Cl^{0}(2,4) \cong \Cl^{0}(4,2) \cong \Mat(4,\C)$. The residue is a soldered invariant, not an algebraic one, and the one metric superselects the tower charge. \emph{Provenance:} \cite[Thms.~3.5 and 4.3, Prop.~4.1]{ButlerT14} (\ref{c:dark-matter.superselected-hence-dark-exactly-stable}, \ref{c:dark-matter.mirror-tower-copy-one}, \ref{c:dark-matter.residue-soldered-invariant-not-algebraic-one}).

\end{proof}

\begin{assembled}[Imported: the conformal completion]\label{c:odd-extensions.imported-conformal-completion}
\changedmark{gray}{reworded}
A shadow of signature $(t,s)$ has conformal completion with ambient of signature $(t+1, s+1)$, one rung up its own tower; for this world the ambient is $(2,4)$ and its conformal spin group acts on the conformal spinor space $\CS \cong \C^{4}$ (twistor space), whose complex unit is $I_6$, through $\Cl^{0}(2,4)$. \emph{Provenance:} \cite[Fact~2.1, Thm.~3.4]{ButlerT9}; \cite[Fact~2.11]{ButlerT12}.
\end{assembled}
\begin{proof}
A shadow of signature $(t,s)$ has conformal completion with ambient of signature $(t+1, s+1)$, one rung up its own tower; for this world the ambient is $(2,4)$ and its conformal spin group acts on the conformal spinor space $\CS \cong \C^{4}$ (twistor space), whose complex unit is $I_6$, through $\Cl^{0}(2,4)$. \emph{Provenance:} \cite[Fact~2.1, Thm.~3.4]{ButlerT9} (\ref{c:correspondence.imported-forced-completion-fibre-quadric}, \ref{c:correspondence.spinor-identification}); \cite[Fact~2.11]{ButlerT12} (\ref{c:standard-model.imported-fock-representation-clifford-algebra-bogoliubov}).

\end{proof}

\begin{fact}[Imported: the frozen point (external ingredients)]\label{c:odd-extensions.imported-frozen-point-external}
The results from outside the series on which \ref{c:odd-extensions.imported-frozen-point} draws, as cited there: S. W. MacDowell and F. Mansouri, Unified geometric theory of gravity and supergravity, Phys. Rev. Lett. 38 (1977) 739.; K. S. Stelle and P. C. West, Spontaneously broken de Sitter symmetry and the gravitational holonomy group, Phys. Rev. D 21 (1980) 1466..
\emph{Source:} S. W. MacDowell and F. Mansouri, Unified geometric theory of gravity and supergravity, Phys. Rev. Lett. 38 (1977) 739.; K. S. Stelle and P. C. West, Spontaneously broken de Sitter symmetry and the gravitational holonomy group, Phys. Rev. D 21 (1980) 1466..
\end{fact}

\begin{assembled}[Imported: the frozen point]\label{c:odd-extensions.imported-frozen-point}
\changedmark{BurntOrange}{changed since the deposit}
On the invertible region of a gravitational field $(\hb, \Omega)$ on flat spacetime (Definition \emph{Gauge Theory Gravity: the gravitational field on flat spacetime}) with $\Lambda \neq 0$, the parallel scale tractor $I$ has squared norm $Q(I) = \Lambda/3$ in the tractor metric of signature $(2,4)$, the square of $\Cl(2,4)$; its orthogonal complement $I^{\perp}$ is a rank-five bundle carrying the induced form of signature $(1,4)$ for $\Lambda > 0$ and $(2,3)$ for $\Lambda < 0$, with stabiliser $SO(1,4)$ or $SO(2,3)$; the Stelle-West vector $V$ is a section of $I^{\perp}$ of squared norm $-3/\Lambda$, its stabiliser is the Lorentz group of $\{I, V\}^{\perp}$, of signature $(1,3)$, and its covariant derivative is the soldering form, the physical frame $e^{\mu} = \hb^{-1}(e^{\mu})$ of the gravitational field, so that in the gauge where $V$ is constant the reduced connection splits as a Lorentz connection, the rotation-gauge field $\Omega$, plus the translation part $e^{\mu}\wedge V/Q(V) = -e^{\mu}\wedge V/\ell^{2}$. \emph{Provenance:} \cite[Fact~2.1, Thm.~3.1]{ButlerT10}; the construction is MacDowell and Mansouri's \cite{MacDowellMansouri1977} in Stelle and West's form \cite{StelleWest1980}.
\end{assembled}
\begin{proof}
Read in Gauge Theory Gravity: the metric $g$ below is the derived metric of the gravitational field $(\hb, \Omega)$ on flat spacetime, the manifold is the invertible region, and the tractor connection is the flat-space $\Cl(2,4)$ gauge field of Theorem \emph{The tractor connection as a rotation-gauge field of the conformal model on flat spacetime}; the index computation below is its reading in the splitting of $g$ (Fact \emph{Imported: the dictionary between Gauge Theory Gravity and curved-spacetime language}).

On the invertible region with $\Lambda \neq 0$ the parallel scale tractor $I$ has squared norm $Q(I) = \Lambda/3$ in the tractor metric of signature $(2,4)$; its orthogonal complement $I^{\perp}$ is a rank-five bundle carrying the induced form of signature $(1,4)$ for $\Lambda > 0$ and $(2,3)$ for $\Lambda < 0$, with stabiliser $SO(1,4)$ or $SO(2,3)$; the Stelle-West vector $V$ is a section of $I^{\perp}$ of squared norm $-3/\Lambda$, its stabiliser is the Lorentz group of $\{I, V\}^{\perp}$, of signature $(1,3)$, and its covariant derivative is the soldering form, so that in the gauge where $V$ is constant the reduced connection splits as a Lorentz connection plus the translation part $e^{\mu}\wedge V/Q(V) = -e^{\mu}\wedge V/\ell^{2}$. \emph{Provenance:} \cite[Fact~2.1, Thm.~3.1]{ButlerT10} (\ref{c:dark-energy.frozen-point-makes-stelle-west-vector}, \ref{c:dark-energy.imported-frozen-point-geometry}); the construction is MacDowell and Mansouri's \cite{MacDowellMansouri1977} in Stelle and West's form \cite{StelleWest1980}.

This restates \ref{c:odd-extensions.imported-frozen-point-external}.

\end{proof}

\section*{What we assume}
The results below are taken as given in this paper; each result that rests on one says so.
\begin{postulate}[P1: the single frozen point]\label{c:dark-matter.p1-single-frozen-point}
The descent was one event freezing one $I$ and making one orientation choice: all towers couple to one metric and are oriented, where they can be, by the one Witt flag of that event (Definition \emph{The Witt flag} [T14, Def 2.10]). The orientation clause is used only on the reversal orbit, where the identification of the flag across the two towers is canonical. \emph{Support:} \cite[Thm.~4.7]{ButlerT6}, \cite[Obs.~7.4]{ButlerT12}; the orientation clause is this paper's. Structural premise.
\end{postulate}

\begin{postulate}[Tameness]\label{c:sheaf.tameness}
The pair $(M, G)$ is a \emph{tame corank-one degeneration} if:
\begin{enumerate}
\item[(T1)] $\Sig$ is a smooth closed hypersurface of $M$;
\item[(T2)] $G$ has corank exactly one at every point of $\Sig$, and the radical lines assemble into a smooth line subbundle $\shL \subset T^*M|_\Sig$, locally spanned by a covector field $k$;
\item[(T3)] (transversality) the degeneration is first-order: the canonical section $\delta$ of Lemma \ref{c:sheaf.invariant-form-t3} below is nowhere vanishing on $\Sig$.
\end{enumerate}
\end{postulate}

\section{What a single generator does}

\begin{definition}[Single extension]\label{c:obj.odd-extensions.single-extension}
The \emph{time extension} of $\Cl(p,q)$ is $\Cl(p+1, q)$ and its \emph{space extension} is $\Cl(p, q+1)$, each with the new generator written $e_{n+1}$ and the old generators, and the subalgebra they generate, called the \emph{shadow}. A single extension is not a transport in the series' sense: the transport adjoins a null pair $e_{+} + e_{-}$ because a zero divisor of vector type is resolved by a nilpotent, and a single generator squaring to $\pm 1$ supplies none \cite[\S 8]{ButlerT1}.
\end{definition}

\begin{lemma}[The volume element in odd dimension]\label{c:odd-extensions.volume-element-odd-dimension}
For $n$ odd, $\omega$ is central in $\Cl(p,q)$ and $\omega^{2} = (-1)^{n(n-1)/2 + q}$. For $n = 5$ this is $\omega^{2} = (-1)^{q}$.
\end{lemma}
\begin{proof}
$e_{i}\omega = (-1)^{n-1}\omega e_{i}$ since $e_{i}$ anticommutes with the $n - 1$ other factors (\ref{c:std-algebra.orthogonal-anticommute}) and commutes with itself; for $n$ odd the sign is $+1$, and the generators generate. Reversing the order of the $n$ factors costs $(-1)^{n(n-1)/2}$ and the contractions $e_{i}^{2}$ cost $(-1)^{q}$. The script verifies centrality and the sign on every diagonal signature of dimension one, three, five and seven.

\end{proof}

\begin{theorem}[Doubling or complexification]\label{c:odd-extensions.doubling-or-complexification}
\changedmark{BrickRed}{corrected since the deposit}
Let $n$ be odd and let $\Cl^{0}$ be the even subalgebra of $\Cl(p,q)$.
\begin{enumerate}
\item[(i)] If $\omega^{2} = +1$, then $\Cl(p,q) \cong \Cl^{0} \oplus \Cl^{0}$, the two summands the ideals of the complementary central idempotents $(1 \pm \omega)/2$, and on the summand of $(1 \pm \omega)/2$ the generator $e_{n}$ acts as $\pm s\,\omega'$, with $s = e_{n}^{2}$ and $\omega' = e_{1}\cdots e_{n-1}$ the volume element of the shadow (so the two summands carry opposite signs).
\item[(ii)] If $\omega^{2} = -1$, then $\Cl(p,q) \cong \Cl^{0} \otimes_{\R} \C$ with $\omega$ as the imaginary unit.
\item[(iii)] $\Cl^{0}(p, q+1) \cong \Cl(p,q)$ and $\Cl^{0}(p+1, q) \cong \Cl(q,p)$.
\end{enumerate}
\end{theorem}
\begin{proof}
(i) With $\omega$ central (Lemma \ref{c:odd-extensions.volume-element-odd-dimension}) and $\omega^{2} = 1$ the elements $P_{\pm} = (1 \pm \omega)/2$ are central, idempotent, orthogonal and sum to one, so the algebra is the direct sum of the ideals $P_{\pm}\Cl$. The map $x \mapsto P_{\pm}x$ restricted to $\Cl^{0}$ is injective, since $\omega$ is odd and $x - \omega x$ can vanish for even $x$ only if $x = 0$, and surjective onto $P_{\pm}\Cl$ because $P_{\pm}y = P_{\pm}(y_{0} \pm \omega y_{1})$ for $y = y_{0} + y_{1}$ split by parity, with $\omega y_{1}$ even; and $P_{\pm}$ is central, so the map is an algebra isomorphism. Finally $\omega = \omega' e_{n}$ with $e_{n}$ commuting with $\omega'$ (it anticommutes with each of the $n - 1$ factors, an even number, \ref{c:std-algebra.orthogonal-anticommute}), so $\omega^{2} = \omega'^{2}e_{n}^{2}$ and $\omega^{2} = 1$ forces $\omega'^{2} = e_{n}^{2} =: s \in \{\pm 1\}$; then $P_{\pm}\omega' = (\omega' \pm s\,e_{n})/2$ and $P_{\pm}e_{n} = (e_{n} \pm s\,\omega')/2 = \pm s\,P_{\pm}\omega'$, so $e_{n}$ acts on the two summands as $+s\,\omega'$ and $-s\,\omega'$. (ii) With $\omega^{2} = -1$ central, $\Cl$ is a $\C$-algebra through $\omega$, and the $\R$-linear map $\Cl^{0} \otimes \C \to \Cl$, $x \otimes (a + bi) \mapsto ax + b\,\omega x$, is an algebra homomorphism, injective by the parity argument, and surjective by dimension ($2 \cdot 2^{n-1} = 2^{n}$, \ref{c:std-algebra.z2-grading}). (iii) For $e_{n}^{2} = -1$ the elements $e_{i}e_{n}$, $i < n$, anticommute pairwise and square to $-e_{i}^{2}e_{n}^{2} = e_{i}^{2}$, so they generate a copy of $\Cl(p,q)$ inside $\Cl^{0}(p, q+1)$, of the full dimension $2^{n-1}$; for $e_{n}^{2} = +1$ they square to $-e_{i}^{2}$ and generate $\Cl(q,p)$. The script verifies (iii), and the identities behind (i) and (ii) (the sign pairing $P_{\pm}e_{n} = \pm s P_{\pm}\omega'$, the dimensions and the spanning, but not the idempotent identities and not the algebra-map properties, which are argued above), over every odd signature of dimensions three and five, and (i) and (ii) in detail on the four extensions of Corollary \ref{c:odd-extensions.four-extensions-oriented-pair}: the idempotents, the two ideals of dimension sixteen, the action of $e_{5}$ as $\pm\omega'$ on them, and the spanning of all thirty-two dimensions by $\Cl(p,q) + \omega\,\Cl(p,q)$ in the complexified cases.

\end{proof}

\begin{corollary}[The four extensions of the oriented pair]\label{c:odd-extensions.four-extensions-oriented-pair}
\changedmark{BurntOrange}{changed since the deposit}
The single extensions of this world and its mirror are, with their residues,
\[
\begin{array}{lllll}
\toprule
\text{extension} & \text{algebra} & \omega^{2} & \text{structure} & r \\
\midrule
\text{this world} + \text{space} & \Cl(1,4) & +1 & \Cl(1,3) \oplus \Cl(1,3) & 5 \\
\text{this world} + \text{time}  & \Cl(2,3) & -1 & \Cl(1,3) \otimes \C \cong \Mat(4,\C), \ \Cl^{0} \cong \Cl(3,1) & 7 \\
\text{mirror} + \text{space}     & \Cl(3,2) & +1 & \Cl(3,1) \oplus \Cl(3,1) & 1 \\
\text{mirror} + \text{time}      & \Cl(4,1) & -1 & \Cl(3,1) \otimes \C \cong \Mat(4,\C), \ \Cl^{0} \cong \Cl(1,3) & 3 \\
\bottomrule
\end{array}
\]
In both towers the generator of the time sign complexifies and the generator of the space sign doubles. The labels follow the uniform convention (time a generator of square $+1$, space one of square $-1$, in both towers); in T14's causal convention the mirror's time is its minus sign, so its space and time extensions are $\Cl(4,1)$ and $\Cl(3,2)$, each the reversal of ours of the same kind, and the convention-free rule is that a generator of square $-1$ doubles and one of square $+1$ complexifies (Proposition \emph{The mirror's extensions under the causal convention}).
\end{corollary}
\begin{proof}
Both shadows have $q$ odd ($3$ and $1$); a time generator leaves $q$ unchanged and a space generator raises it by one, so by Lemma \ref{c:odd-extensions.volume-element-odd-dimension} the time extensions have $\omega^{2} = -1$ and the space extensions $\omega^{2} = +1$, and Theorem \ref{c:odd-extensions.doubling-or-complexification} gives the structures. The even subalgebras are from Theorem \ref{c:odd-extensions.doubling-or-complexification}(iii): $\Cl^{0}(1,4) \cong \Cl(1,3)$, $\Cl^{0}(2,3) \cong \Cl(3,1)$, $\Cl^{0}(3,2) \cong \Cl(3,1)$, $\Cl^{0}(4,1) \cong \Cl(1,3)$. The residues are $p - q \bmod 8$. The identification $\Cl(1,3) \otimes \C \cong \Mat(4,\C)$ (\ref{c:std-algebra.sta-matrix}) is Fact \ref{c:odd-extensions.imported-representatives-dimension-nearest-four-types} complexified.

\end{proof}

\begin{remark}[the added axis is the pseudoscalar]\label{c:note.odd-extensions.added-axis-pseudoscalar}
In the doubled case the fifth generator is nothing new: on each summand it acts as the pseudoscalar $\omega' = e_{1}e_{2}e_{3}e_{4}$ of the shadow, up to sign, so the extra space direction of $\Cl(1,4)$ is the chirality operator of $\Cl(1,3)$ promoted from a $4$-blade to a vector, with the choice of sign the only datum the extension adds. In the complexified case the fifth generator is $e_{5} = \pm\,\omega\,\omega'^{-1}$, that is $i\gamma_{5}$ in the physicists' notation, and the algebra is the complexified Dirac algebra.
\end{remark}

\section{Where the mirror sits}

\begin{theorem}[The mirror inside the complexified extension, and absent from the doubled one]\label{c:odd-extensions.mirror-inside-complexified-extension-absent-doubled}
\changedmark{gray}{reworded}
In $\Cl(2,3)$, the time extension of this world with $e_{5}^{2} = +1$:
\begin{enumerate}
\item[(i)] the generators $e_{1}, \ldots, e_{4}$ generate a copy of $\Cl(1,3)$, this world's shadow;
\item[(ii)] the bivectors $e_{\mu}e_{5}$, $\mu = 1, \ldots, 4$, anticommute pairwise, square to $(-1, +1, +1, +1)$, and generate the whole even subalgebra, which is a copy of $\Cl(3,1)$, the mirror's shadow;
\item[(iii)] the two copies meet exactly in the even part of the shadow, $\Cl^{0}(1,3) \cong \Mat(2,\C)$, of dimension eight, the Lorentz algebra with its Weyl spinors; and together with $\omega$ they span everything, $\Cl(2,3) = \Cl(1,3) \oplus \omega\,\Cl(1,3)$;
\item[(iv)] no automorphism of $\Cl(2,3)$ carries the copy of $\Cl(1,3)$ onto the copy of $\Cl(3,1)$; the natural map between their generating sets is multiplication by the central $\omega$, $e_{\mu} \mapsto \omega e_{\mu}$, which squares to $-e_{\mu}^{2}$ (the Wick rotation) and is complex-linear, not an automorphism of the real algebra.
\end{enumerate}
In $\Cl(1,4)$, the space extension with $e_{5}^{2} = -1$: the generators $e_{\mu}$ and the bivectors $e_{\mu}e_{5}$ both generate copies of $\Cl(1,3)$, the elements $\omega e_{\mu}$ generate a third, and $\Cl(3,1)$ admits no unital embedding into $\Cl(1,4)$ at all. The same statements hold for the mirror's two extensions with the roles of $\Cl(1,3)$ and $\Cl(3,1)$ exchanged.
\end{theorem}
\begin{proof}
(i) is the definition of the shadow. (ii) is Theorem \ref{c:odd-extensions.doubling-or-complexification}(iii) with $e_{5}^{2} = +1$: the products square to $-e_{\mu}^{2}$, which is $(-1, +1, +1, +1)$, and generate a subalgebra of dimension sixteen inside the even subalgebra of dimension sixteen. (iii) The shadow's even part consists of even elements of $\Cl(2,3)$, so it lies in the intersection; the shadow's odd part consists of odd elements of $\Cl(2,3)$, since the generators $e_{1}, \ldots, e_{4}$ are odd there, so it meets the even subalgebra in zero; the intersection is exactly $\Cl^{0}(1,3)$, of dimension eight, and $\Cl^{0}(1,3) \cong \Cl(3,0)$ (\ref{c:std-algebra.sta-matrix}) $\cong \Mat(2,\C)$ (\ref{c:std-algebra.pauli-algebra}) is the Pauli algebra. For the spanning statement: $\omega$ times a blade of the shadow is, up to sign, the complementary blade of the shadow times $e_{5}$, so $\Cl(1,3) + \omega\,\Cl(1,3)$ contains every blade of $\Cl(2,3)$ (the script checks the dimension); this is the shadow's version of Theorem \ref{c:odd-extensions.doubling-or-complexification}(ii), which decomposes by the even subalgebra instead. (iv) An automorphism carries a subalgebra to an isomorphic subalgebra, and $\Cl(1,3) \not\cong \Cl(3,1)$ by Fact \ref{c:odd-extensions.imported-representatives-dimension-nearest-four-types}. With $\omega$ central and $\omega^{2} = -1$, $(\omega e_{\mu})^{2} = -e_{\mu}^{2}$ and the $\omega e_{\mu}$ anticommute pairwise, so they generate a $\Cl(3,1)$; this is the complex-linear identification of the two real forms and it is not an automorphism of the real algebra, since $\omega e_{\mu}$ is a $4$-blade, not a vector.

For $\Cl(1,4)$: Theorem \ref{c:odd-extensions.doubling-or-complexification}(iii) with $e_{5}^{2} = -1$ gives $(e_{\mu}e_{5})^{2} = e_{\mu}^{2}$, and $\omega^{2} = +1$ gives $(\omega e_{\mu})^{2} = e_{\mu}^{2}$, so all three sets generate $\Cl(1,3)$. For the non-embedding: $\Cl(1,4) \cong A \oplus A$ with $A = \Cl(1,3) \cong \Mat(2,\HH)$ by Theorem \ref{c:odd-extensions.doubling-or-complexification}(i) and Fact \ref{c:odd-extensions.imported-representatives-dimension-nearest-four-types}. A unital homomorphism $\phi: \Mat(4,\R) \to A \oplus A$ composed with the projection to a summand on which $\phi(1) \neq 0$ is a unital homomorphism $\Mat(4,\R) \to A$; its kernel is a two-sided ideal of the simple algebra $\Mat(4,\R)$ not containing $1$, hence zero; so it is injective between real algebras of equal dimension sixteen, hence an isomorphism, contradicting $\Mat(4,\R) \not\cong \Mat(2,\HH)$ (Fact \ref{c:odd-extensions.imported-representatives-dimension-nearest-four-types}). The statements for the mirror's extensions follow by exchanging the roles of $p$ and $q$ throughout. The script verifies the squares, the anticommutation, the dimensions of the generated subalgebras and of their intersections with the even subalgebra, and the Wick-rotated squares, for all four extensions.

\end{proof}

\begin{proposition}[The translation bivectors of the frozen point]\label{c:odd-extensions.translation-bivectors-frozen-point}
\changedmark{gray}{reworded}
Under Fact \ref{c:odd-extensions.imported-frozen-point}, write $\Cl(I^{\perp})$ for the Clifford bundle of the rank-five bundle $I^{\perp}$, and let $e_{5} = V/|V|$ be the unit Stelle-West direction, so that the Lorentz frame $e_{1}, \ldots, e_{4}$ of $\{I,V\}^{\perp}$ and $e_{5}$ together generate $\Cl(I^{\perp})$ fibrewise. Then:
\begin{enumerate}
\item[(i)] for $\Lambda > 0$, $\Cl(I^{\perp}) \cong \Cl(1,4)$, $e_{5}^{2} = -1$, and the translation bivectors $e_{\mu}e_{5}$ of the MacDowell-Mansouri connection generate this world's shadow $\Cl(1,3)$;
\item[(ii)] for $\Lambda < 0$, $\Cl(I^{\perp}) \cong \Cl(2,3)$, $e_{5}^{2} = +1$, and the translation bivectors generate a copy of the mirror's abstract algebra $\Cl(3,1)$, while this world's shadow is the subalgebra generated by the Lorentz frame (what this copy is and is not: Proposition \ref{c:odd-extensions.belongs-tower} and Corollary \ref{c:odd-extensions.two-readings-which-which}).
\end{enumerate}
In either case the fibrewise algebra $\Cl(I^{\perp})$ is one of the two extensions of Corollary \ref{c:odd-extensions.four-extensions-oriented-pair}, and the Poincar\'e contraction $\ell \to \infty$ leaves the Clifford type of the translation algebra unchanged.
\end{proposition}
\begin{proof}
Fact \ref{c:odd-extensions.imported-frozen-point} gives the signature of $I^{\perp}$ as $(1,4)$ or $(2,3)$ by the sign of $\Lambda$, the sign of $Q(V) = -3/\Lambda$, and $\{I,V\}^{\perp}$ of signature $(1,3)$ in both cases; so $e_{5}^{2} = \operatorname{sign}Q(V)$ is $-1$ for $\Lambda > 0$ and $+1$ for $\Lambda < 0$, and $\Cl(I^{\perp})$ is the space extension or the time extension of $\Cl(1,3)$ respectively. The translation part of the reduced connection is $e^{\mu}\wedge V/\ell$, whose Clifford image is the bivector $e_{\mu}e_{5}$ up to the scalar $1/\ell$; the algebra these generate is Theorem \ref{c:odd-extensions.mirror-inside-complexified-extension-absent-doubled}: $\Cl(1,3)$ in the first case, $\Cl(3,1)$ in the second. The contraction rescales the generators by $1/\ell$ and rescaling does not change the signs of the squares.

\end{proof}

\begin{remark}[what this says, and what it does not]\label{c:note.odd-extensions.what-says-what-it-does-not}
The de Sitter algebra of this world is two copies of this world, and the anti-de Sitter algebra of this world is this world merged with its mirror, with the mirror sitting on the momenta. This is a statement about the fibrewise Clifford algebra of $I^{\perp}$, which Fact \ref{c:odd-extensions.imported-frozen-point} supplies; it is not a statement that a negative cosmological constant would turn matter into mirror matter. The tower charge is soldered and superselected (Fact \ref{c:odd-extensions.imported-copy-theorem-soldered-residue}), and Proposition \ref{c:odd-extensions.every-path-mirror-crosses-two-time} below shows that no path of sign changes at fixed dimension four connects the two oriented towers without passing through a sterile signature. What the proposition does establish is that the sign of $\Lambda$ selects which of two non-isomorphic algebras the frozen point's translation generators live in, and that only one of them contains the mirror. The observed cosmological constant is positive \cite{Planck2018}, so the algebra realised in this universe is the doubled $\Cl(1,4)$ with no copy of the mirror in it; the mirror on the momenta is a statement about the other sign. The same source gives the present Hubble rate as exceeding the asymptotic rate $\sqrt{\Lambda/3}$ of Theorem \ref{c:odd-extensions.exceptional-algebra-residue-completion-both-sides}(ii) by the factor $\Omega_{\Lambda}^{-1/2}$.
\end{remark}

\section{The rung's even subalgebra is the anti-de Sitter algebra}

\begin{theorem}[The even conformal algebra is $\Cl(2,3)$]\label{c:odd-extensions.twistor-algebra}
\changedmark{gray}{reworded}
Let $\Cl(2,4)$ be the ambient of this world's conformal completion (Fact \ref{c:odd-extensions.imported-conformal-completion}). Then:
\begin{enumerate}
\item[(i)] $\Cl^{0}(2,4) \cong \Cl(2,3) \cong \Cl(4,1) \cong \Mat(4,\C)$, the two odd algebras being generated inside the even subalgebra by $e_{i}e_{6}$ for a spacelike $e_{6}$ and by $e_{i}e_{1}$ for a timelike $e_{1}$ respectively; the rung's volume element $\omega_{6}$ (the conformal pseudoscalar, $I_6$ up to sign) is even, squares to $-1$ and commutes with every even element, and is the complex unit of $\Mat(4,\C)$;
\item[(ii)] inside $\Cl^{0}(2,4)$ both oriented shadows occur: the bivectors $e_{i}e_{6}$ for $i = 2, \ldots, 5$ generate a $\Cl(1,3)$ and the bivectors $e_{i}e_{1}$ for $i = 2, \ldots, 5$ generate a $\Cl(3,1)$;
\item[(iii)] every input to the copy theorem (Fact \ref{c:odd-extensions.imported-copy-theorem-soldered-residue}) is an object of $\Cl^{0}(2,4)$ or of its complexified modules, hence of the complexification in which $\Cl(1,3)$ and $\Cl(3,1)$ are two real forms of one algebra;
\item[(iv)] signature reversal is not an isomorphism of the odd extensions: $\Cl(2,3) \not\cong \Cl(3,2)$ and $\Cl(1,4) \not\cong \Cl(4,1)$, the centre being $\C$ in $\Cl(2,3)$ and $\Cl(4,1)$ and $\R \oplus \R$ in $\Cl(3,2)$ and $\Cl(1,4)$.
\end{enumerate}
\end{theorem}
\begin{proof}
(i) Theorem \ref{c:odd-extensions.doubling-or-complexification}(iii) applied to $\Cl(2,4)$ with a spacelike distinguished generator gives $\Cl^{0}(2,4) \cong \Cl(2,3)$, and with a timelike one $\Cl^{0}(2,4) \cong \Cl(4,1)$; both are $\Mat(4,\C)$ by Corollary \ref{c:odd-extensions.four-extensions-oriented-pair}. For $n = 6$ even, $\omega_{6}$ anticommutes with every vector, hence commutes with every even element, and $\omega_{6}^{2} = (-1)^{15}(-1)^{4} = -1$ (\ref{c:std-algebra.pseudoscalar-commutation}, \ref{c:primer.sign}); being even it lies in $\Cl^{0}(2,4)$ and is central there, and since $\Cl^{0}(2,4) \cong \Mat(4,\C)$ has centre $\C$, $\omega_{6}$ is the complex unit up to sign. (ii) The squares are $(e_{i}e_{6})^{2} = -e_{i}^{2}e_{6}^{2} = e_{i}^{2} = (1, -1, -1, -1)$ for $i = 2, \ldots, 5$ with the signature $(+,+,-,-,-,-)$, and $(e_{i}e_{1})^{2} = -e_{i}^{2} = (-1, 1, 1, 1)$; both sets anticommute pairwise and generate subalgebras of dimension sixteen. (iii) The conformal spin group is generated by the bivectors of $\Cl(2,4)$, which are even, and acts on $\CS$ through $\Cl^{0}(2,4)$ (Fact \ref{c:odd-extensions.imported-conformal-completion}); the conformal spinor space $\CS \cong \C^{4}$ (twistor space) is the irreducible module of $\Cl^{0}(2,4) \cong \Mat(4,\C)$; the isotropic structure is a structure on $V_{\C} = \C^{6}$, the complexification of the ambient; and the mass formula is a bivector norm in the even subalgebra. (iv) An isomorphism of algebras restricts to an isomorphism of centres; the centre of an odd Clifford algebra is spanned by $1$ and $\omega$: a basis blade commutes with $e_{i}$ exactly when $e_{i}$ anticommutes with an even number of its factors, which for a blade of grade $k$ containing $e_{i}$ requires $k - 1$ even and for one not containing it requires $k$ even, and only the scalar and the volume element satisfy both conditions for every $i$; since $x \mapsto e_{i}xe_{i}^{-1}$ is diagonal in the blade basis, an element commutes with every $e_{i}$ exactly when each of its blades does, so the centre is spanned by those two (the script confirms dimension two over every odd signature of dimension at most five, and dimension one over every even signature of dimension at most six), and is $\R \oplus \R$ when $\omega^{2} = +1$ and $\C$ when $\omega^{2} = -1$; Corollary \ref{c:odd-extensions.four-extensions-oriented-pair} gives the signs. The script verifies (i), (ii) and the centre dimensions.

\end{proof}

\begin{corollary}[The copy theorem's blind spot]\label{c:odd-extensions.copy-theorems-blind-spot}
\changedmark{BurntOrange}{changed since the deposit}
The copy theorem identifies this world and its mirror by working in an algebra, $\Cl^{0}(2,4) \cong \Cl(2,3)$, that contains both real forms and is the complexification of each, its complex unit being the conformal pseudoscalar $I_6$ that spans its centre with $1$. It is therefore exactly blind to the distinction that a single extension sees: the two oriented towers have isomorphic conformal spin groups, isomorphic conformal spinor modules (twistor modules) and one complexified ambient, but odd extensions that signature reversal does not identify (Theorem \ref{c:odd-extensions.twistor-algebra}(iv)), whose centres differ. The shadows themselves are non-isomorphic (Fact \ref{c:odd-extensions.imported-representatives-dimension-nearest-four-types}); every object the copy theorem is built from, the even subalgebra of the rung, the conformal spin group and the conformal spinor module $\CS$ (Theorem \ref{c:odd-extensions.twistor-algebra}(iii)), identifies the two towers, while the rungs themselves, $\Cl(2,4)$ and $\Cl(4,2)$, do not, being of the two towers' types by Facts \ref{c:odd-extensions.imported-towers-clock} and \ref{c:odd-extensions.imported-representatives-dimension-nearest-four-types}; the odd extension is the first construction above the shadow, by dimension, that keeps them apart, and it does so through an element, the bare fifth generator, that no derived quantity of the series uses (Proposition \ref{c:odd-extensions.two-towers-one-curved-universe}).
\end{corollary}
\begin{proof}
Theorem \ref{c:odd-extensions.twistor-algebra}(iii) and (iv). The remaining clauses rest on Fact \ref{c:odd-extensions.imported-representatives-dimension-nearest-four-types}, which supplies the non-isomorphic shadow types $\Cl(1,3) \cong \Mat(2,\HH)$ and $\Cl(3,1) \cong \Mat(4,\R)$; on Fact \ref{c:odd-extensions.imported-towers-clock}, whose $(1,1)$ ascent preserves the residue and so gives the rungs the two towers' types; and on Proposition \ref{c:odd-extensions.two-towers-one-curved-universe}, which supplies that the element distinguishing the two towers' intermediate algebras is one no derived quantity of the series uses.

\end{proof}

\begin{remark}[the holographic coincidence]\label{c:note.odd-extensions.holographic-coincidence}
Writing, as here, $\Cl(2,4)$ for the Clifford algebra of the flat space in which $\mathrm{AdS}_{5}$ is the hyperboloid, and likewise for the other hyperboloids: the rung is the algebra of $\mathrm{AdS}_{5}$, whose isometry group $SO(2,4)$ is the conformal group of this world; the intermediate $\Cl(2,3)$ is the algebra of $\mathrm{AdS}_{4}$ and, by Proposition \ref{c:odd-extensions.four-extensions-are-four-odd-towers}, the conformal ambient of the three-dimensional residue-$7$ shadow, the group theory of $\mathrm{AdS}_{4}/\mathrm{CFT}_{3}$; the de Sitter intermediate $\Cl(1,4)$ is the ambient of the Euclidean three-dimensional shadow $(0,3)$, also of residue $5$, the group theory of $\mathrm{dS}/\mathrm{CFT}$. All of this is the standard fact that the isometry group of $\mathrm{AdS}_{d+1}$ is the conformal group in $d$ dimensions, which any conformal field theory in $d$ dimensions shares whether or not a bulk exists; the note adds only the residues. The coincidence does expose a difference in what two frameworks say about one pair of algebras. Holography takes the bulk and its boundary to be one physics in two descriptions. The series takes the residue of a sector to be the signature of its soldered space (Proposition \ref{c:odd-extensions.belongs-tower}), which a conformal class carries as well as a metric, so the boundary of $\mathrm{AdS}_{4}$ has residue $7$ and lies, by Fact \ref{c:odd-extensions.imported-towers-clock}, in a tower distinct from the bulk's; whether a conformal class with no chosen scale is also a superselected sector, Fact \ref{c:odd-extensions.imported-copy-theorem-soldered-residue} speaking of the one metric, is a separate question and is not settled here; and the series' superselection is a statement about operators soldered to one metric, which a holographic dictionary is not. Whether that last difference means the two frameworks are speaking about different things, or that one of them is wrong about this pair, is not argued anywhere in the series and is not argued here; the remark records the question and claims nothing about the dynamical content of holography, to which the series has no counterpart.
\end{remark}

\section{The clock, and the paths between the oriented towers}

\begin{proposition}[The four extensions are the four odd towers]\label{c:odd-extensions.four-extensions-are-four-odd-towers}
\changedmark{gray}{reworded}
The residues of the four extensions are $5$, $7$, $1$, $3$ (Corollary \ref{c:odd-extensions.four-extensions-oriented-pair}), and they land on the odd towers in two roles: $\Cl(1,4)$ and $\Cl(4,1)$, this world's space extension and the mirror's time extension, are the representatives at the dimension nearest four of the residues $5$ and $3$ (Fact \ref{c:odd-extensions.imported-representatives-dimension-nearest-four-types}), the two five-dimensional dark towers; $\Cl(2,3)$ and $\Cl(3,2)$, this world's time extension and the mirror's space extension, are the conformal ambients (Fact \ref{c:odd-extensions.imported-conformal-completion}) of the representatives $(1,2)$ and $(2,1)$ of the residues $7$ and $1$, the two three-dimensional dark towers. This is a statement about abstract algebras: the frozen point's algebra $\Cl(I^{\perp})$ of Proposition \ref{c:odd-extensions.translation-bivectors-frozen-point} coincides with a five-dimensional shadow for $\Lambda > 0$ and with a three-dimensional ambient for $\Lambda < 0$ as an algebra, and Proposition \ref{c:odd-extensions.belongs-tower} says that as a soldered object it is neither.
\end{proposition}
\begin{proof}
$(1,4)$ and $(4,1)$, the space and time extensions of Corollary \ref{c:odd-extensions.four-extensions-oriented-pair}, are listed in Fact \ref{c:odd-extensions.imported-representatives-dimension-nearest-four-types} at $r = 5$ and $r = 3$. The conformal completion (Fact \ref{c:odd-extensions.imported-conformal-completion}) of $(1,2)$ is $(2,3)$ and of $(2,1)$ is $(3,2)$, residues $7$ and $1$, preserved by the $(1,1)$ step (Fact \ref{c:odd-extensions.imported-towers-clock}). The last sentence is Proposition \ref{c:odd-extensions.translation-bivectors-frozen-point} read against the first two, with Proposition \ref{c:odd-extensions.belongs-tower} for the soldering. The script checks every residue.

\end{proof}

\begin{proposition}[Every path to the mirror crosses the two-time fork]\label{c:odd-extensions.every-path-mirror-crosses-two-time}
At fixed dimension four, a change of the sign of one generator moves the signature $(p, 4-p)$ to $(p \pm 1, 4 - p \mp 1)$ and the residue by $\pm 2$. Every sequence of such changes from $(1,3)$ to $(3,1)$ passes through $(2,2)$, residue $0$, the two-time sterile fork; the other neighbour of $(1,3)$, the Euclidean $(0,4)$ at residue $4$, is a dead end from which the only continuation returns to $(1,3)$.
\end{proposition}
\begin{proof}
The signatures at dimension four are the five points $p = 0, \ldots, 4$ on a line, a sign change moves $p$ by one, and a walk from $p = 1$ to $p = 3$ visits $p = 2$; $p = 0$ has one neighbour. The script enumerates every non-revisiting walk of length at most six and finds exactly one, $(1,3) \to (2,2) \to (3,1)$.

\end{proof}

\begin{remark}[no mechanism]\label{c:note.odd-extensions.no-mechanism}
Together with Fact \ref{c:odd-extensions.imported-copy-theorem-soldered-residue} this closes the question of a physical passage between the oriented towers within the series: the tower charge is superselected by the one metric, the only algebra in which both towers are present identifies them by the complex unit and not by any automorphism (Theorem \ref{c:odd-extensions.mirror-inside-complexified-extension-absent-doubled}(iv)), and the only continuous route through signatures at dimension four runs through a signature with no causal structure. The two towers share gravity, and by Proposition \ref{c:odd-extensions.translation-bivectors-frozen-point} the sign of $\Lambda$; that is the whole of the shared channel.
\end{remark}

\section{The frozen point splits the ascent's pair}

\phantomsection\label{c:note.odd-extensions.setting-frozen-point-splits-ascents-pair}
The rung of this world is $\Cl(2,4)$, obtained from $\Cl(1,3)$ by adjoining a hyperbolic pair, one generator of each sign, whose null combinations are the zero divisors the transport resolves (Fact \ref{c:odd-extensions.imported-towers-clock}, Definition \ref{c:obj.odd-extensions.single-extension}). The frozen point of the dark-energy paper is a vector in that rung, and its norm is the cosmological constant. This section shows that $I$ and the Stelle-West vector $V$ are the ascent's pair, signed by $\Lambda$, and that the odd extensions of the section ``What a single generator does'' are the intermediate algebras of this world's own ascent, with the degenerate case at $\Lambda = 0$ the wall of the sheaf paper.

\begin{assembled}[Imported: the frozen point at $\Lambda = 0$, and the degenerate algebras]\label{c:odd-extensions.imported-frozen-point-degenerate-algebras}
\changedmark{gray}{reworded}
At $\Lambda = 0$ the parallel scale tractor is null, its stabiliser is parabolic (the Poincar\'e-type stabiliser of a null direction), and the rotation to the reductive de Sitter stabiliser happens exactly when $Q(I) \neq 0$; the identity $Q(I)\,Q(V) = -1$ holds for $\Lambda \neq 0$. A quadratic space with a one-dimensional radical has Clifford algebra $\Cl(p,q;1)$ with a nilpotent generator $n$, $n^{2} = 0$, and $n\,\Cl(p,q;1)$ is a two-sided ideal of square zero; the series treats these algebras, their nilpotent ideals and their one-sided completions in the sheaf paper, and a nilpotent generator is the vector-type zero divisor whose resolution is the transport. \emph{Provenance:} \cite[Fact~2.1, Thm.~3.1]{ButlerT10}; \cite[Thm.~3.1, Prop.~3.4, Thm.~4.4]{ButlerT4}; \cite[\S 4]{ButlerT1}.
\end{assembled}
\begin{proof}
At $\Lambda = 0$ the parallel scale tractor is null, its stabiliser is parabolic (the Poincar\'e-type stabiliser of a null direction), and the rotation to the reductive de Sitter stabiliser happens exactly when $Q(I) \neq 0$; the identity $Q(I)\,Q(V) = -1$ holds for $\Lambda \neq 0$. A quadratic space with a one-dimensional radical has Clifford algebra $\Cl(p,q;1)$ with a nilpotent generator $n$, $n^{2} = 0$, and $n\,\Cl(p,q;1)$ is a two-sided ideal of square zero; the series treats these algebras, their nilpotent ideals and their one-sided completions in the sheaf paper, and a nilpotent generator is the vector-type zero divisor whose resolution is the transport. \emph{Provenance:} \cite[Fact~2.1, Thm.~3.1]{ButlerT10} (\ref{c:dark-energy.frozen-point-makes-stelle-west-vector}, \ref{c:dark-energy.imported-frozen-point-geometry}); \cite[Thm.~3.1, Prop.~3.4, Thm.~4.4]{ButlerT4} (\ref{c:sheaf.degeneration-sheaf}, \ref{c:sheaf.resolution-ideal-theoretically}, \ref{c:sheaf.one-sided-completion}); \cite[\S 4]{ButlerT1}.

\end{proof}

\begin{theorem}[$I$ and $V$ are the ascent's pair, signed by $\Lambda$]\label{c:odd-extensions.are-ascents-pair-signed}
\changedmark{gray}{reworded}
In the tractor bundle of signature $(2,4)$ let $I$ be the frozen point, $Q(I) = \Lambda/3$, and for $\Lambda \neq 0$ let $V$ be the Stelle-West vector. Then:
\begin{enumerate}
\item[(i)] $I$ and $V$ span a hyperbolic plane: $\ip{I}{V} = 0$, $Q(I)$ and $Q(V)$ have opposite signs, and the unit directions $e_{I}$, $e_{V}$ anticommute with each other and with the Lorentz frame of $\{I,V\}^{\perp}$;
\item[(ii)] $I^{\perp}$ is spanned by the Lorentz frame and $e_{V}$, and $\Cl(I^{\perp})$ is the subalgebra of the rung generated by these five: the space extension $\Cl(1,4)$ for $\Lambda > 0$ ($e_{V}^{2} = -1$) and the time extension $\Cl(2,3)$ for $\Lambda < 0$ ($e_{V}^{2} = +1$); adjoining $e_{I}$ returns the whole rung, so the ascent factors as
\[
\Cl(1,3) \;\subset\; \Cl(I^{\perp}) \;\subset\; \Cl(2,4),
\]
with the sign of $\Lambda$ deciding which of the pair's two signed directions is adjoined first;
\item[(iii)] at $\Lambda = 0$ the frozen point, written $I_{0}$ from here on, is null and lies in its own complement, $I^{\perp}$ is degenerate of signature $(1,3;1)$, and $\Cl(I^{\perp}) \cong \Cl(1,3;1)$ is the degenerate algebra with $I$ itself the nilpotent generator: $I^{2} = 0$, $I$ anticommutes with the Lorentz frame, and $I\,\Cl(1,3;1)$ is a sixteen-dimensional ideal of square zero;
\item[(iv)] the Lorentz frame $\{I,V\}^{\perp}$ has signature $(1,3)$ for every value of $\Lambda$, including zero.
\end{enumerate}
\end{theorem}
\begin{proof}
Write the rung's hyperbolic plane orthogonal to the Lorentz frame as $\mathrm{span}(e_{+}, e_{-})$ with $e_{+}^{2} = +1$, $e_{-}^{2} = -1$; the frozen point is $I = c\,e_{+} + d\,e_{-}$ with $c^{2} - d^{2} = \Lambda/3$, and the orthogonal vector in the plane is $w = d\,e_{+} + c\,e_{-}$ with $Q(w) = d^{2} - c^{2} = -Q(I)$ and $\ip{I}{w} = cd - cd = 0$. (i) For $\Lambda \neq 0$, $V$ is the section of $I^{\perp}$ of norm $-3/\Lambda$ orthogonal to the Lorentz frame (Fact \ref{c:odd-extensions.imported-frozen-point}), so $V$ is a multiple of $w$ and $Q(I)\,Q(V) < 0$ is the identity of Fact \ref{c:odd-extensions.imported-frozen-point-degenerate-algebras}. Orthogonal vectors anticommute in a Clifford algebra. (ii) $I^{\perp} = \{I,V\}^{\perp} \oplus \mathrm{span}(w)$ by dimension, and $e_{V}^{2} = \operatorname{sign}Q(V) = -\operatorname{sign}\Lambda$; so the five generators have signature $(1,3) + (0,1) = (1,4)$ for $\Lambda > 0$ and $(1,3) + (1,0) = (2,3)$ for $\Lambda < 0$, which are the two extensions of Corollary \ref{c:odd-extensions.four-extensions-oriented-pair}. The five together with $e_{I}$ are an orthogonal basis of the rung's vector space, so they generate $\Cl(2,4)$. (iii) At $c = d$, $Q(I) = 0$ and $w = I$, so $I \in I^{\perp}$, the induced form has the one-dimensional radical $\mathrm{span}(I)$, and the Clifford algebra of the five generators has one null generator: $\Cl(1,3;1)$, with the ideal and its square as in Fact \ref{c:odd-extensions.imported-frozen-point-degenerate-algebras}. (iv) The Lorentz frame is orthogonal to both $e_{+}$ and $e_{-}$ and its signature does not depend on $c, d$. The script verifies (i) to (iv) on exact rational instances with $Q(I) = 3$, $-3$ and $0$: the orthogonality relations, the signatures $(1,4)$, $(2,3)$ and $(1,3;1)$ by exact congruence, the anticommutation, the generation of all sixty-four dimensions when $Q(I) \neq 0$, the nilpotent generator and its ideal in $\Cl(1,3;1)$, and the signature $(1,3)$ of the Lorentz frame in all three cases.

\end{proof}

\begin{corollary}[The cosmological constant deforms the wall]\label{c:odd-extensions.cosmological-constant-deforms-wall}
The intermediate algebra of this world's ascent is a one-parameter family in $\Lambda$: the degenerate $\Cl(1,3;1)$ at $\Lambda = 0$, whose nilpotent generator is the zero divisor the transport is built to resolve, deformed to the doubled $\Cl(1,4)$ for $\Lambda > 0$ and the complexified $\Cl(2,3)$ for $\Lambda < 0$. The sign of $\Lambda$ fixes the sign given to the ascent's null generator (the generator's square is $-\Lambda/3$, the opposite sign), and the two odd extensions are the two ways of giving it one.
\end{corollary}
\begin{proof}
Theorem \ref{c:odd-extensions.are-ascents-pair-signed}(ii) and (iii); the family is $I = c\,e_{+} + d\,e_{-}$ with $c^{2} - d^{2} = \Lambda/3$.

\end{proof}

\begin{remark}[what this adds to the dark-energy paper]\label{c:note.odd-extensions.what-adds-dark-energy-paper}
The dark-energy paper proves the parabolic-to-reductive rotation of the stabiliser at $\Lambda \neq 0$ at the level of groups. Theorem \ref{c:odd-extensions.are-ascents-pair-signed} is the same rotation at the level of the Clifford algebra: the null direction of the parabolic case is the nilpotent generator of a degenerate algebra, and the reductive case is that generator given a sign. That the tangent image keeps signature $(1,3)$ whatever $\Lambda$ does is Theorem \ref{c:odd-extensions.are-ascents-pair-signed}(iv) read through Fact \ref{c:odd-extensions.imported-copy-theorem-soldered-residue}.
\end{remark}

\section{The soldering decides}

\phantomsection\label{c:note.odd-extensions.setting-soldering-decides}
the section ``The clock, and the paths between the oriented towers'' placed the four extensions on the four odd towers as algebras. The dark-matter paper's definition of the residue is not algebraic, and applying it settles which tower the frozen point's algebra belongs to.

\begin{proposition}[$\Cl(I^{\perp})$ belongs to this tower]\label{c:odd-extensions.belongs-tower}
Under the dark-matter paper's Theorem 3.5 the residue of a sector is the signature of the soldered vector space, the image of the tangent space in the algebra with $v^{2} = g(v,v)$ (Fact \ref{c:odd-extensions.imported-copy-theorem-soldered-residue}). For the algebra $\Cl(I^{\perp})$ of Theorem \ref{c:odd-extensions.are-ascents-pair-signed}:
\begin{enumerate}
\item[(i)] the soldered vector space is the Lorentz frame $\{I,V\}^{\perp} = e(TM)$, the image of the soldering form, of signature $(1,3)$ for every $\Lambda$ (Theorem \ref{c:odd-extensions.are-ascents-pair-signed}(iv)); $V$ is not in it, being orthogonal to $e(TM)$ with $\nabla V = e$ (Fact \ref{c:odd-extensions.imported-frozen-point});
\item[(ii)] hence $\Cl(I^{\perp})$ has tower charge $6$: it is a sub-object of this world's rung, the intermediate algebra of this world's own ascent, and not the ground shadow of the residue-$5$ or residue-$7$ tower whose abstract algebra it shares;
\item[(iii)] the coincidence $\Cl(I^{\perp}) \cong \Cl(1,4)$ for $\Lambda > 0$ is induced by an isometry of quadratic spaces ($I^{\perp}$ has signature $(1,4)$ outright) but by no map of soldered data: the isometry carries $e_{V}$, which is not soldered, to a fifth tangent direction of the five-dimensional shadow, which is, and carries the four-dimensional soldered space into a proper subspace of that shadow's tangent space.
\end{enumerate}
\end{proposition}
\begin{proof}
(i) is Fact \ref{c:odd-extensions.imported-frozen-point} with Theorem \ref{c:odd-extensions.are-ascents-pair-signed}(iv). (ii) The residue is the signature of the soldered space by definition (Fact \ref{c:odd-extensions.imported-copy-theorem-soldered-residue}), and that signature is $(1,3)$, residue $6$; the algebra generated by the soldered space together with the unsoldered $e_{V}$ lies inside the rung generated by the soldered space together with the ascent's pair. (iii) A five-dimensional shadow of residue $5$ is, by the same definition, an algebra whose soldered space has signature $(1,4)$, so all five generators are tangent; under the one four-dimensional metric of the dark-matter paper's Premise P1 (\ref{c:dark-matter.p1-single-frozen-point}) this needs a fifth tangent direction beyond the four, an internal direction in the sense of that paper's projection lemma \cite[Lem.~6.1]{ButlerT14} (\ref{c:dark-matter.projection-every-towers-content-dilutes-ones}). The isometry $I^{\perp} \cong \R^{1,4}$ identifies $e_{V}$ with that direction, so it does not carry the soldered space onto the soldered space, which is what the dark-matter paper's Theorem 3.5(i) requires of a map that preserves the residue.

\end{proof}

\begin{corollary}[The two readings, and which is which]\label{c:odd-extensions.two-readings-which-which}
The two readings of $\Cl(1,4)$ that occur in the series are two solderings of one algebra, and the series' definition of the residue assigns them to different towers. The fifth direction of the frozen point's algebra is not a spacetime direction and never was: it is the Stelle-West vector, the normal to the de Sitter hyperboloid, with fixed norm $\ell = \sqrt{3/\Lambda}$, and its algebra is this tower's. The five-dimensional dark towers, if populated, are soldered to a fifth tangent direction that the frozen point does not supply, whose size $L$ remains the datum the series does not fix. Proposition \ref{c:odd-extensions.four-extensions-are-four-odd-towers} therefore holds as a statement about algebras, and only about algebras.
\end{corollary}
\begin{proof}
Proposition \ref{c:odd-extensions.belongs-tower}.

\end{proof}

\begin{corollary}[$\Lambda$ is not a coupling between towers]\label{c:odd-extensions.not-coupling-between-towers}
\changedmark{gray}{reworded}
The cosmological constant is a datum of this tower's ascent (Theorem \ref{c:odd-extensions.are-ascents-pair-signed}), and its coincidence with the residue-$5$ tower's abstract algebra carries no structure across the superselection, since the dark-matter paper's exactness argument is precisely that isomorphisms between towers that do not respect the soldering carry nothing (Fact \ref{c:odd-extensions.imported-copy-theorem-soldered-residue}). The coupling rule the dark-matter paper's odd towers need is not supplied by $\Lambda$.
\end{corollary}
\begin{proof}
Proposition \ref{c:odd-extensions.belongs-tower}(iii) and Fact \ref{c:odd-extensions.imported-copy-theorem-soldered-residue}. The first sentence rests on Theorem \ref{c:odd-extensions.are-ascents-pair-signed}, which supplies $\Lambda$ as a datum of this tower's ascent, the sign of $Q(I) = \Lambda/3$ on the ascent's pair.

\end{proof}

\begin{remark}[what the copies are]\label{c:note.odd-extensions.what-copies-are}
the section ``What a single generator does'' to the section ``The rung's even subalgebra is the anti-de Sitter algebra'' are about this world's own ascent. The copy of the mirror's algebra that Proposition \ref{c:odd-extensions.translation-bivectors-frozen-point}(ii) finds on the translation bivectors for $\Lambda < 0$ is not the mirror tower, whose vectors are soldered; it is the reason the two oriented towers' intermediate algebras are non-isomorphic (Corollary \ref{c:odd-extensions.copy-theorems-blind-spot}), and that non-isomorphism is a statement about the two towers' own ascents under their own solderings, unaffected by this section. Likewise the odd towers' algebras that Proposition \ref{c:odd-extensions.four-extensions-are-four-odd-towers} finds among the four extensions are their abstract algebras only.
\end{remark}

\section{Running the transport on the wall}

\phantomsection\label{c:note.odd-extensions.setting-running-transport-wall}
the section ``The frozen point splits the ascent's pair'' put the frozen point at $\Lambda = 0$ on the wall: the nilpotent generator of $\Cl(1,3;1)$. The series has a procedure for walls, the transport, and this section runs it and asks what it fixes and what it leaves free. The answer is that it fixes the rung and a canonical null frame in the rung's new plane, and nothing about the frozen point's position in that plane beyond what its norm, the cosmological constant, records.

\begin{assembled}[Imported: the minimal completion of a degenerate algebra, and Rule 1]\label{c:odd-extensions.imported-minimal-completion-degenerate-algebra-rule}
A real quadratic space with nondegenerate part of signature $(p,q)$ and radical of dimension $r$ admits a nondegenerate completion; every completion has dimension at least $p + q + 2r$; the minimal ones have signature exactly $(p+r, q+r)$; and any two minimal completions are isometric by an isometry fixing the original space pointwise, so the minimal ambient Clifford algebra $\Cl(p+r, q+r)$ is unique up to an isomorphism fixing the embedded $\Cl(p,q;r)$. In the model completion the radical is paired with its dual, $Q(\rho, \varphi) = 1$, $Q(\varphi,\varphi) = 0$, $\varphi$ orthogonal to the nondegenerate part. For a type-N obstruction in a nondegenerate algebra the increment $(+1,+1)$ is forced outright; for a type-C obstruction the completed algebra is unique but the sign of the adjoined generator's square is a geometric datum the algebra forgets. \emph{Provenance:} \cite[Thm.~2.3 and its model]{ButlerT2}; \cite[Thm.~4.13, Thm.~5.7 and Rem.~5.8]{ButlerT1}.
\end{assembled}
\begin{proof}
A real quadratic space with nondegenerate part of signature $(p,q)$ and radical of dimension $r$ admits a nondegenerate completion; every completion has dimension at least $p + q + 2r$; the minimal ones have signature exactly $(p+r, q+r)$; and any two minimal completions are isometric by an isometry fixing the original space pointwise, so the minimal ambient Clifford algebra $\Cl(p+r, q+r)$ is unique up to an isomorphism fixing the embedded $\Cl(p,q;r)$. In the model completion the radical is paired with its dual, $Q(\rho, \varphi) = 1$, $Q(\varphi,\varphi) = 0$, $\varphi$ orthogonal to the nondegenerate part. For a type-N obstruction in a nondegenerate algebra the increment $(+1,+1)$ is forced outright; for a type-C obstruction the completed algebra is unique but the sign of the adjoined generator's square is a geometric datum the algebra forgets. \emph{Provenance:} \cite[Thm.~2.3 and its model]{ButlerT2} (\ref{c:conformal.minimal-nondegenerate-completion}); \cite[Thm.~4.13, Thm.~5.7 and Rem.~5.8]{ButlerT1} (\ref{c:transport.minimality}, \ref{c:transport.one-generator-minimal-minimal-completions-classified}).

\end{proof}

\begin{theorem}[The transport constrains the frozen point only through its norm]\label{c:odd-extensions.transport-constrains-frozen-point-only-through}
Let $W = I_{0}^{\perp}$, the degenerate space of signature $(1,3;1)$ at $\Lambda = 0$ with radical spanned by the null frozen point $I_{0}$, written $I$ in this theorem, and let $\widetilde W$ be its minimal completion. Then:
\begin{enumerate}
\item[(i)] $\widetilde W$ has signature $(2,4)$ and is the rung; the isometry of Fact \ref{c:odd-extensions.imported-minimal-completion-degenerate-algebra-rule} fixing $W$ pointwise is unique, so the null partner $\varphi$ of $I$ with $Q(I,\varphi) = 1$, $Q(\varphi,\varphi) = 0$, $\varphi \perp \{I\}^{\perp} \cap W$ is canonical: the transport supplies the hyperbolic plane $\mathrm{span}(I, \varphi)$ together with a distinguished null basis of it;
\item[(ii)] every vector orthogonal to the Lorentz frame is $aI + b\varphi$ with norm $2ab$; a frozen point of cosmological constant $\Lambda$ is such a vector with $2ab = \Lambda/3$, and the sign of $\Lambda$ is the sign of the product of its components along the null generator and its canonical partner;
\item[(iii)] the isometries of the completion that fix the Lorentz frame act on the plane as $O(1,1)$, whose orbits on nonzero vectors are exactly the level sets of the norm; so the only invariant of the frozen point under the symmetries the transport supplies is its norm $\Lambda/3$: the cosmological constant, sign and magnitude, is the one modulus of the frozen point that the transport leaves free;
\item[(iv)] no symmetry of the completion carries a frozen point of positive norm to one of negative norm, and the two intermediate algebras $\Cl(I_{\Lambda}^{\perp})$ for the two signs are non-isomorphic (Theorem \ref{c:odd-extensions.twistor-algebra}(iv)); the two signs are inequivalent under every symmetry the construction has, and the choice between them is decided by nothing in the transport.
\end{enumerate}
\end{theorem}
\begin{proof}
(i) Fact \ref{c:odd-extensions.imported-minimal-completion-degenerate-algebra-rule} with $(p,q,r) = (1,3,1)$ gives $(2,4)$. An isometry of $\widetilde W$ fixing $W$ pointwise sends $\varphi$ to a vector $x$ with $Q(x, w) = Q(\varphi, w)$ for all $w \in W$ and $Q(x,x) = 0$; the five linear conditions have rank five in six unknowns, so $x = \varphi + cI$, and $Q(x,x) = 2c\,Q(I,\varphi) = 2c$ forces $c = 0$. (ii) $Q(aI + b\varphi, aI + b\varphi) = 2ab\,Q(I,\varphi) = 2ab$; the vector is orthogonal to the Lorentz frame because $I$ and $\varphi$ are. (iii) The plane $\mathrm{span}(I,\varphi)$ is the orthogonal complement of the Lorentz frame, so an isometry fixing the frame restricts to an isometry of the plane, and the isometry group of a hyperbolic plane is $O(1,1)$: the boosts $(a,b) \mapsto (ta, b/t)$, $t > 0$, the swap $(a,b) \mapsto (b,a)$, and their products with $-1$. The boosts are transitive on each branch of each hyperbola $2ab = \text{const} \neq 0$ and on each of the four null rays; the element $-1$ exchanges the two branches of a hyperbola, and the swap exchanges the two null axes; so the orbits are the level sets of the norm. (iv) Isometries preserve the norm; the non-isomorphism is Theorem \ref{c:odd-extensions.twistor-algebra}(iv). The script verifies the signature, the rank-five linear system and its unique null solution, the norm formula on instances of both signs and zero, the boost and swap invariances, and the centres.

\end{proof}

\begin{corollary}[The cosmological constant is undetermined within the series]\label{c:odd-extensions.cosmological-constant-undetermined-within-series}
The transport, run on the wall, returns the rung and a canonical null frame in its new plane, and stops; the frozen point's position in that plane is data the transport receives, not data it produces, and the one invariant of that position is the norm $\Lambda/3$, sign and magnitude together. Within the series the sign is a choice between two options that the construction makes inequivalent (Theorem \ref{c:odd-extensions.transport-constrains-frozen-point-only-through}(iv)) and that nothing in the construction decides, which is the status the dark-matter paper's Premise P2 gives to the choice of which oriented tower is the visible one; the analogy is one of status, since P2's choice is binary and this one is a real modulus. Nothing here says the sign has no cause: only that the transport is not it.
\end{corollary}
\begin{proof}
Theorem \ref{c:odd-extensions.transport-constrains-frozen-point-only-through}(iii) and (iv), with \cite[Premise P2]{ButlerT14} for the status.

\end{proof}

\begin{remark}[the dictionary with the ascent's null pair]\label{c:note.odd-extensions.dictionary-ascents-null-pair}
\changedmark{gray}{reworded}
The correspondence between the frozen point and the pair the transport adjoins, stated in the root paper's normalisation (which fixes the null basis up to the dilation, one hyperbolic plane being isometric to any other). The adjoined plane has the signed basis $e_{\pm}$ and the null basis $v = \tfrac12(e_{+} + e_{-})$, $v^{*} = \tfrac12(e_{+} - e_{-})$, with $v^{2} = v^{*2} = 0$ and $vv^{*} + v^{*}v = 1$, the pair the root paper adjoins and the tower paper reads as the one fermionic mode of the rung \cite[\S 2.2]{ButlerT1}, \cite[Prop.~4.1, Thm.~4.2]{ButlerT3}. The note's canonical frame is the same two null rays, $v = I_{0}$ and $v^{*} = \varphi/2$. Then: at $\Lambda = 0$ the frozen point is the null generator $v$ itself, the zero divisor of vector type, and the direction the completion adjoins is its partner $v^{*}$, canonical by Theorem \ref{c:odd-extensions.transport-constrains-frozen-point-only-through}(i); the Stelle-West vector does not exist there. For $\Lambda \neq 0$ the frozen point is a combination of both null generators,
\[
I_{\Lambda} \;=\; \alpha\, v + \beta\, v^{*}, \qquad Q(I_{\Lambda}) = \alpha\beta = \Lambda/3 ,
\]
and the Stelle-West vector is the same two components with the relative sign reversed,
\[
V \;\propto\; \alpha\, v - \beta\, v^{*}, \qquad Q(V) \propto -\alpha\beta ,
\]
orthogonal to $I_{\Lambda}$ and to the Lorentz frame. So $I$ and $V$ are the sum and difference of one pair of null components, which is why they form the signed basis up to scale and a boost (Theorem \ref{c:odd-extensions.are-ascents-pair-signed}(i)); the sign of $\Lambda$ is the relative sign of $\alpha$ and $\beta$, de Sitter when they agree and anti-de Sitter when they differ, and the wall is either of them vanishing. The dilation acts as $v \mapsto tv$, $v^{*} \mapsto v^{*}/t$, an automorphism of the relation $vv^{*} + v^{*}v = 1$, and $\alpha\beta$ is its invariant (Proposition \ref{c:odd-extensions.boost-dilation-ratio-gauge}); in the frame of Theorem \ref{c:odd-extensions.transport-constrains-frozen-point-only-through}, $\alpha = a$ and $\beta = 2b$. In the tower paper's reading of $v$ and $v^{*}$ as the annihilation and creation operators of the rung's mode, the frozen point of a $\Lambda = 0$ universe is a pure annihilator and a nonzero $\Lambda$ is a creation component, the product of the two coefficients being $\Lambda/3$; this is what that reading says verbatim; only the vanishing of the creation component is invariant, its size being the dilation gauge, and the note attaches no claim to the reading. The script verifies the null pair, its normalisation, the identification with the canonical frame, and the norms and orthogonality of $\alpha v \pm \beta v^{*}$ on instances of both signs and zero.
\end{remark}

\begin{remark}[two bits, compared]\label{c:note.odd-extensions.two-bits-compared}
The root paper's Rule 2 leaves a sign free too, the square of the generator adjoined to a type-C obstruction, and records that the completed algebra forgets it. The sign left free here is of a different kind. The completed algebra, the rung $\Cl(2,4)$, forgets it in the same way; but the intermediate algebra of Theorem \ref{c:odd-extensions.are-ascents-pair-signed} remembers it, the doubled $\Cl(1,4)$ against the complexified $\Cl(2,3)$, because the base $\Cl(1,3)$ is even-dimensional and exactly one of the two single extensions changes the parity of $q$ (Lemma \ref{c:odd-extensions.volume-element-odd-dimension}). So the sign of $\Lambda$ is a geometric datum with an algebraic witness one level below the rung, and the witness is what the section ``The rung's even subalgebra is the anti-de Sitter algebra'' used to separate the oriented towers.
\end{remark}

\begin{remark}[what the transport stance does and does not say about $\Lambda$]\label{c:note.odd-extensions.what-transport-stance-does-does-not}
The programme reads a zero divisor as a signpost rather than a prohibition, and Theorem \ref{c:odd-extensions.are-ascents-pair-signed}(iii) puts $\Lambda = 0$ on the wall. It is tempting to read the stance as saying that a universe does not sit on the wall, hence $\Lambda \neq 0$. Theorem \ref{c:odd-extensions.transport-constrains-frozen-point-only-through} is what the stance actually delivers: the transport off the wall is the passage from $\Cl(1,3;1)$ to the rung, which is the same rung for every $\Lambda$, and it says nothing about where in the rung's new plane the frozen point lies. The dark-energy paper's finding that the norm is a parameter of the laws under the action as written is the dynamical form of the same statement, and this note does not improve on it.
\end{remark}

\section{Two negative results}

\phantomsection\label{c:note.odd-extensions.setting-two-negative-results}
Theorem \ref{c:odd-extensions.transport-constrains-frozen-point-only-through}(iii) says the frozen point's two components in the canonical null frame are related by a symmetry of the completion. This section adds that the symmetry is a gauge symmetry of the dark-energy paper's action, so the components are unphysical and not merely symmetry-related; and it shows that the sign of $\Lambda$, though it acts differently on the two oriented towers, does not select which of them is the visible one.

\begin{proposition}[The boost is the dilation; the ratio $b/a$ is gauge]\label{c:odd-extensions.boost-dilation-ratio-gauge}
\changedmark{gray}{reworded}
In the rung, with the canonical null frame realised as $I = (e_{+} + e_{-})/2$, $\varphi = e_{+} - e_{-}$ in the hyperbolic plane orthogonal to the Lorentz frame:
\begin{enumerate}
\item[(i)] the bivector $B = I \wedge \varphi = -e_{+}e_{-}$ squares to $+1$, commutes with every Lorentz bivector and with every Lorentz vector, and anticommutes with $e_{\pm}$: it is the generator of the dilations of the conformal model;
\item[(ii)] the rotor $R = \cosh\theta + B\sinh\theta$ fixes the Lorentz frame and acts on the plane by $I \mapsto tI$, $\varphi \mapsto \varphi/t$ with $t = e^{2\theta}$, so the boosts of Theorem \ref{c:odd-extensions.transport-constrains-frozen-point-only-through}(iii) are the dilations;
\item[(iii)] hence the ratio $b/a$ of the frozen point's components is a dilation gauge and only the product $2ab = \Lambda/3$ is invariant, and by Fact \ref{c:dark-energy.point-infinity-has-no-physical-fluctuation} the gauge is not fixed by the dynamics either: the components are unphysical, not merely symmetry-related, and the frame the transport supplies carries no quantity beyond $\Lambda$ itself.
\end{enumerate}
\end{proposition}
\begin{proof}
(i) $I\varphi = \tfrac12(e_{+} + e_{-})(e_{+} - e_{-}) = 1 - e_{+}e_{-}$, so $I \wedge \varphi = -e_{+}e_{-}$ and $B^{2} = -e_{+}^{2}e_{-}^{2} = +1$; a bivector of a plane commutes with vectors orthogonal to the plane and anticommutes with vectors in it (\ref{c:std-algebra.orthogonal-anticommute}). (ii) $R = \cosh\theta + B\sinh\theta$ is the exponential of the bivector $\theta B$ (\ref{c:std-algebra.bivector-exponential}), and $R\tilde R = \cosh^{2}\theta - \sinh^{2}\theta = 1$, $R$ fixes what commutes with $B$, and on the plane $RxR^{-1}$ is the hyperbolic rotation by $2\theta$, whose eigenvectors are the null directions with eigenvalues $e^{\pm 2\theta}$. (iii) is (ii) with Theorem \ref{c:odd-extensions.transport-constrains-frozen-point-only-through}(iii) and Fact \ref{c:dark-energy.point-infinity-has-no-physical-fluctuation}. The script verifies (i) and (ii) with the rational rotor $\cosh\theta = 5/4$, $\sinh\theta = 3/4$ ($t = 4$): unit, fixing $e_{1}, \ldots, e_{4}$, scaling $I$ by $4$ and $\varphi$ by $1/4$, and carrying a frozen point along its hyperbola with $2ab$ preserved.

\end{proof}

\begin{proposition}[The two towers on one gravitational field]\label{c:odd-extensions.two-towers-one-curved-universe}
\changedmark{gray}{reworded}
Both oriented towers are soldered to the one gravitational field $(\hb, \Omega)$ on flat spacetime (Definition \emph{Gauge Theory Gravity: the gravitational field on flat spacetime}), the mirror reading its derived metric with the sign reversed, $-g$, so the mirror's tractor bundle has signature $(4,2)$ (this world's tractors being the vectors of $\Cl(2,4)$ carried by the tractor gauge field of Theorem \emph{The tractor connection as a rotation-gauge field of the conformal model on flat spacetime}, the mirror's those of $\Cl(4,2)$) and its frozen point has norm $-Q(I)$. Then:
\begin{enumerate}
\item[(i)] in a de Sitter universe the intermediate algebra of this world's ascent is the doubled $\Cl(1,4)$ and the mirror's is the complexified $\Cl(4,1)$; in an anti-de Sitter universe they are the complexified $\Cl(2,3)$ and the doubled $\Cl(3,2)$. The rule is: doubled if and only if (de Sitter and this world) or (anti-de Sitter and mirror);
\item[(ii)] the even subalgebras $\Cl^{0}(1,4)$ and $\Cl^{0}(4,1)$ have the same presentation, the bivectors $e_{\mu}e_{5}$ anticommuting with squares $(+,-,-,-)$, so the two towers' conformal-to-de-Sitter spin groups, spinor modules and MacDowell-Mansouri bivector algebras coincide; the derived content of the dark-matter paper's copy theorem, which runs through the spin group, does not see the doubling;
\item[(iii)] what differs is the bare fifth generator alone: $e_{5}^{2} = -1$ with $e_{5} = \pm\omega'$ on the two summands in the doubled case, against $e_{5}^{2} = +1$ with $e_{5} = \omega'^{-1}\omega$, a central complex unit times the pseudoscalar, in the complexified case.
\end{enumerate}
Consequently the sign of $\Lambda$ does not select the visible tower: it distinguishes the two towers' intermediate algebras, but through an element no derived quantity of the series uses, and that element is the unit vector along $V$, which Proposition \ref{c:odd-extensions.belongs-tower} places outside the soldered space and Fact \ref{c:dark-energy.point-infinity-has-no-physical-fluctuation} places in the pure-gauge sector.
\end{proposition}
\begin{proof}
The metric $g$ below is the derived metric $g(a,b) = \hb^{-1}(a) \cdot \hb^{-1}(b)$ of the gravitational field and $-g$ the mirror's reading of it; its Levi-Civita connection, Ricci tensor and scalar curvature are read from $\Omega$ and $\mathcal{R}$ by Fact \emph{Imported: the dictionary between Gauge Theory Gravity and curved-spacetime language}(ii) and (iv). (i) The tractor bundle of a conformal class of signature $(p,q)$ carries a metric of signature $(p+1,q+1)$ (Fact \ref{c:odd-extensions.imported-frozen-point}), so the class $[-g]$ of signature $(3,1)$ has a tractor metric of signature $(4,2)$; in a scale $g$ the tractor metric is $2\sigma\rho + g^{-1}(\mu,\mu)$ and the scale tractor of $\sigma$ is $(\sigma, \nabla\sigma, -(\Delta\sigma + J\sigma)/n)$ with $J$ a multiple of the scalar curvature; under $g \mapsto -g$ the Levi-Civita connection and the Ricci tensor are unchanged while the inverse metric, the Laplacian and the scalar curvature change sign, so $\rho \mapsto -\rho$; the map $\Phi: (\sigma, \mu, \rho) \mapsto (\sigma, \mu, -\rho)$ from the tractor bundle of $[g]$ to that of $[-g]$ is scale-independent, intertwines the two tractor connections, carries the scale tractor of $g$ to that of $-g$, and satisfies $h_{-g}(\Phi X, \Phi X) = -h_{g}(X, X)$. Hence the mirror's frozen point has norm $-Q(I)$ and its $I^{\perp}$ has the reversed signature of ours: $(4,1)$ against $(1,4)$ and $(3,2)$ against $(2,3)$; the doubling or complexification is Corollary \ref{c:odd-extensions.four-extensions-oriented-pair}. (ii) Theorem \ref{c:odd-extensions.doubling-or-complexification}(iii): with $e_{5}^{2} = -1$ in $\Cl(1,4)$ and $e_{5}^{2} = +1$ in $\Cl(4,1)$ the products $e_{\mu}e_{5}$ square to $e_{\mu}^{2}$ and $-e_{\mu}^{2}$ respectively, which are both $(+,-,-,-)$ since the two Lorentz frames are reversed; a Clifford algebra is determined by such a presentation. That the copy theorem's inputs are spin-group objects is Fact \ref{c:odd-extensions.imported-copy-theorem-soldered-residue}. (iii) is Theorem \ref{c:odd-extensions.doubling-or-complexification}(i) and Remark \ref{c:note.odd-extensions.added-axis-pseudoscalar}. In the closing clause, the unit vector along $V$ lies outside the soldered space by Proposition \ref{c:odd-extensions.belongs-tower}(i) and in the pure-gauge sector by Fact \ref{c:dark-energy.point-infinity-has-no-physical-fluctuation}. The script verifies the signatures of the mirror's $I^{\perp}$ on the same rational instances, the two-by-two rule from the signs of $\omega^{2}$, the common presentation of the even subalgebras, and the squares and factorisations of $e_{5}$.

\end{proof}

\begin{remark}[the one operator that sees the difference]\label{c:note.odd-extensions.one-operator-sees-difference}
An operator linear in $e_{5}$ alone, a term $\bar\psi e_{5}\psi$, is the one place the doubling would show: $e_{5}$ is the real pseudoscalar in one tower and the pseudoscalar times a central complex unit in the other, which in the physicists' notation is $\gamma_{5}$ against $i\gamma_{5}$, an anti-Hermitian against a Hermitian bilinear. But $e_{5}$ is the unit Stelle-West direction, unsoldered by Proposition \ref{c:odd-extensions.belongs-tower}(i) and pure gauge by Fact \ref{c:dark-energy.point-infinity-has-no-physical-fluctuation}; so either a coupling to it can be physical, in which case Fact \ref{c:dark-energy.point-infinity-has-no-physical-fluctuation} is being read too strongly in Proposition \ref{c:odd-extensions.boost-dilation-ratio-gauge}(iii), or it cannot, in which case the mirror sector is observationally identical to this one at every level the series reaches. The note takes the second reading, consistently with Proposition \ref{c:odd-extensions.boost-dilation-ratio-gauge}, and records the operator as the one place the first reading would show; the Hermiticity statement depends on a Dirac-representation convention the note does not fix.
\end{remark}

\section{The wall in the cosmological constant}

\phantomsection\label{c:note.odd-extensions.setting-wall-cosmological-constant}
Theorem \ref{c:odd-extensions.are-ascents-pair-signed} gives a one-parameter family of intermediate algebras, nondegenerate for $\Lambda \neq 0$ and degenerate at $\Lambda = 0$. The sheaf paper has a theory of exactly such families: tame corank-one degenerations, their exceptional algebras, the residue of the metric at the wall, and a conservation law that says the degeneracy of a wall in spacetime can be displaced but never destroyed. This section runs that theory on the $\Lambda$ line. The family is a tame wall in the sheaf paper's sense, and the rung completes it on both sides, which the conservation law forbids for the completions it covers; the family lies in the case the conservation theorem's proof excludes by hypothesis, a complement whose signature switches across the wall, and that case is generic to any corank-one family of hyperplanes in a fixed nondegenerate space. The section's content is the identification of the dark-energy paper's frozen-point geometry with the sheaf paper's tame structure, and the placement of the $\Lambda$ family in that excluded case.

\begin{assembled}[Imported: tame walls, the exceptional algebra, the residue, and conservation]\label{c:odd-extensions.imported-tame-walls-exceptional-algebra-residue}
A pair $(M,G)$, $G$ a smooth symmetric form on a bundle over $M$, is a tame corank-one degeneration if (T1) the degenerate locus $\Sigma$ is a smooth closed hypersurface, (T2) the corank is exactly one on $\Sigma$ with the radical lines a smooth line bundle spanned locally by $k$, and (T3) the dying eigenvalue crosses zero at nonzero speed along every transversal, invariantly: the section $\delta(k \otimes k) = dg_{k}|_{\Sigma}$ of $\operatorname{Hom}(L^{\otimes 2}, N^{*}\Sigma)$ is nowhere zero. The signatures on the two sides differ by one unit, $(p,q)$ where the dying eigenvalue is positive against $(p-1, q+1)$. Rescaling the dying direction by $t^{-1/2}$ along a transversal gives a continuous family of nondegenerate algebras through the wall whose special fibre, the exceptional algebra, is the Clifford algebra of the surviving form orthogonal to the rate of degeneration, of signature $(p,q)$ on the side approached. The inverse form has a canonical simple pole at the wall whose residue is $\delta^{-1}$, giving $N^{*}\Sigma \cong L^{\otimes 2}$: the dying line is a square root of the conormal bundle and the adjoined partner line a square root of the normal bundle. Conservation, absolute form: for a tame wall with both sides nonempty, any symmetric form on any bundle containing the original, standard outside a neighbourhood of the wall in the sense of being $G \perp K$ there with $K$ nondegenerate of the \emph{same} signature on both exterior regions, degenerates somewhere; the constancy of $\operatorname{sig} K$ across the wall is necessary, a completion allowed to switch $K$'s signature being able to match the two sides. \emph{Provenance:} \cite[Hyp.~2.1, Lem.~2.2, Rem.~2.3, Prop.~6.1, Thm.~6.2, Thm.~5.1]{ButlerT4}.
\end{assembled}
\begin{proof}
A pair $(M,G)$, $G$ a smooth symmetric form on a bundle over $M$, is a tame (\ref{c:sheaf.tameness}) corank-one degeneration if (T1) the degenerate locus $\Sigma$ is a smooth closed hypersurface, (T2) the corank is exactly one on $\Sigma$ with the radical lines a smooth line bundle spanned locally by $k$, and (T3) the dying eigenvalue crosses zero at nonzero speed along every transversal, invariantly: the section $\delta(k \otimes k) = dg_{k}|_{\Sigma}$ of $\operatorname{Hom}(L^{\otimes 2}, N^{*}\Sigma)$ is nowhere zero. The signatures on the two sides differ by one unit, $(p,q)$ where the dying eigenvalue is positive against $(p-1, q+1)$. Rescaling the dying direction by $t^{-1/2}$ along a transversal gives a continuous family of nondegenerate algebras through the wall whose special fibre, the exceptional algebra, is the Clifford algebra of the surviving form orthogonal to the rate of degeneration, of signature $(p,q)$ on the side approached. The inverse form has a canonical simple pole at the wall whose residue is $\delta^{-1}$, giving $N^{*}\Sigma \cong L^{\otimes 2}$: the dying line is a square root of the conormal bundle and the adjoined partner line a square root of the normal bundle. Conservation, absolute form: for a tame wall with both sides nonempty, any symmetric form on any bundle containing the original, standard outside a neighbourhood of the wall in the sense of being $G \perp K$ there with $K$ nondegenerate of the \emph{same} signature on both exterior regions, degenerates somewhere; the constancy of $\operatorname{sig} K$ across the wall is necessary, a completion allowed to switch $K$'s signature being able to match the two sides. \emph{Provenance:} \cite[Hyp.~2.1, Lem.~2.2, Rem.~2.3, Prop.~6.1, Thm.~6.2, Thm.~5.1]{ButlerT4} (\ref{c:sheaf.metrics-residue-square-root-twist}, \ref{c:sheaf.invariant-form-t3}, \ref{c:sheaf.exceptional-algebra-deformation-normal-cone}, \ref{c:sheaf.conservation-degeneracy-absolute-form}).

\end{proof}

\begin{proposition}[The $\Lambda$ line is a tame wall]\label{c:odd-extensions.line-tame-wall}
\changedmark{gray}{reworded}
Fix the dilation gauge by a smooth nowhere-vanishing $a(\Lambda)$ and write the frozen point in the canonical null frame as $I_{\Lambda} = a I_{0} + b\varphi$ with $2ab = \Lambda/3$ (Theorem \ref{c:odd-extensions.transport-constrains-frozen-point-only-through}). Let $w_{\Lambda} = aI_{0} - b\varphi$, so that $I_{\Lambda}^{\perp}$ is spanned by the Lorentz frame and $w_{\Lambda}$. Then the family $(I_{\Lambda}^{\perp}, \text{induced form})$ over the $\Lambda$ line is a tame corank-one degeneration in the sense of Fact \ref{c:odd-extensions.imported-tame-walls-exceptional-algebra-residue}:
\begin{enumerate}
\item[(i)] its Gram matrix in the frame $(e_{1}, \ldots, e_{4}, w_{\Lambda})$ is $\operatorname{diag}(1,-1,-1,-1,-\Lambda/3)$ for every $\Lambda$ and every gauge, so the degenerate locus is the point $\Lambda = 0$ (T1), the radical there is the line spanned by $I_{0}$ (T2), and the dying eigenvalue $-\Lambda/3$ crosses zero with nonzero speed (T3), $\delta(w_{0} \otimes w_{0}) = -\tfrac13\,d\Lambda$ in the parametrisation by $\Lambda$, the $3$ being that of $Q(I) = \Lambda/3$;
\item[(ii)] the signatures on the two sides are $(2,3)$ for $\Lambda < 0$ and $(1,4)$ for $\Lambda > 0$: the wall separates the anti-de Sitter and de Sitter intermediate algebras, and crossing it changes the sign of one direction's square, which is the sheaf paper's signature-change wall in a parameter rather than in spacetime;
\item[(iii)] the gauge $a = b$, which makes $I_{\Lambda}$ proportional to $I_{0} + \varphi$, is not smooth at the wall and is not defined for $\Lambda < 0$; the tame structure needs a smooth gauge, and the scalar invariants (the dying eigenvalue, its crossing speed, the signatures) do not depend on which, while the identification of the radical vector with $I_{0}$ at a fixed normalisation does, $w_{0} = a(0)\,I_{0}$.
\end{enumerate}
\end{proposition}
\begin{proof}
(i) With $2ab = \Lambda/3$ (Theorem \ref{c:odd-extensions.transport-constrains-frozen-point-only-through}(ii)), $Q(w_{\Lambda}) = -2ab = -\Lambda/3$, $\ip{w_{\Lambda}}{I_{\Lambda}} = ab - ab = 0$, and both are orthogonal to the Lorentz frame; the Gram is block diagonal because $I_{0}$ and $\varphi$ are. The radical at $\Lambda = 0$ is spanned by $w_{0} = aI_{0}$. The crossing speed is $d(-\Lambda/3)/d\Lambda = -1/3 \neq 0$, so (T1) to (T3) of Fact \ref{c:odd-extensions.imported-tame-walls-exceptional-algebra-residue} hold. (ii) is the sign of $-\Lambda/3$ read against Corollary \ref{c:odd-extensions.four-extensions-oriented-pair}. (iii) $a = b = \sqrt{\Lambda/6}$. The script verifies the Gram, the corank at $\Lambda = 0$ only, the gauge independence of the dying eigenvalue, the exact slope $-1/3$ at six nonzero rational values of $\Lambda$, and the two signatures.

\end{proof}

\begin{theorem}[Exceptional algebra, residue, and the completion on both sides]\label{c:odd-extensions.exceptional-algebra-residue-completion-both-sides}
\changedmark{gray}{reworded}
For the tame wall of Proposition \ref{c:odd-extensions.line-tame-wall}:
\begin{enumerate}
\item[(i)] the sheaf paper's rescaling of the dying direction, $w_{\Lambda} \mapsto w_{\Lambda}/\sqrt{|\Lambda|}$, is proportional, by the fixed factor $\sqrt 3$, to the dark-energy paper's normalisation of the Stelle-West vector to unit length, $e_{V} = V/|V|$, both being normalisations of the same line; the exceptional algebra at the wall is the Lorentz frame orthogonal to the rate of degeneration, $\Cl(1,3) \perp \langle +\tfrac13\rangle \cong \Cl(2,3)$ approached from $\Lambda < 0$ and $\Cl(1,3) \perp \langle -\tfrac13\rangle \cong \Cl(1,4)$ from $\Lambda > 0$, the cross terms vanishing identically so that the rescaled family is exactly, not asymptotically, block diagonal;
\item[(ii)] the residue of the inverse form at the wall is $-3\,w_{0} \otimes w_{0}$, the coefficient fixed by the normalisation $Q(w_{\Lambda}) = -\Lambda/3$: the pole is the Stelle-West norm, $\Lambda\,Q(V) = -3$, one algebraic step from the imported identity of Fact \ref{c:odd-extensions.imported-frozen-point-degenerate-algebras}; the sheaf paper's identification $N^{*}\Sigma \cong L^{\otimes 2}$ reads $d\Lambda \cong -3\, w_{0} \otimes w_{0}$, which is $-3\,I_{0} \otimes I_{0}$ when $a(0)^{2} = 1$, so the dying direction carries the weight of $|\Lambda|^{1/2}$ and the weightless normalisation of the frozen point's complement is by $\sqrt{|\Lambda|/3}$, which for $\Lambda > 0$ is the asymptotic de Sitter rate $H = \sqrt{\Lambda/3}$, the series' $H^{2} = Q(I)$;
\item[(iii)] the rung, with its tractor metric of constant signature $(2,4)$, is a nondegenerate completion of the family on both sides of the wall, of rank one more, the completing direction at the wall being the partner $\varphi$ and away from it the frozen point $I_{\Lambda}$ itself; this is not particular to the cosmological constant: for any $(p,q)$ the family of hyperplanes $\mathrm{span}(f_{1}, \ldots, f_{n}, u + \tfrac{t}{2}v)$ of $\R^{p,q} \perp H$, $H$ hyperbolic with null basis $u, v$, $\ip{u}{v} = 1$, is a tame corank-one family with dying eigenvalue $t$ completed on both sides by the fixed $(p+1,q+1)$, and the $\Lambda$ family is this model with $(p,q) = (1,3)$, $u = I_{0}$, $v = \varphi$, $t = -\Lambda/3$;
\item[(iv)] this does not contradict the conservation of degeneracy, and lies in the case its proof excludes by hypothesis: the complement of $I_{\Lambda}^{\perp}$ in the rung is $\mathrm{span}(I_{\Lambda})$ with $K = Q(I_{\Lambda}) = \Lambda/3$, whose signature is $(1,0)$ on one side and $(0,1)$ on the other, so the completion is not standard in the sense of Fact \ref{c:odd-extensions.imported-tame-walls-exceptional-algebra-residue}, and the totals $(2,3) + (0,1) = (1,4) + (1,0) = (2,4)$ match because the complement's line changes type across the wall; on the $\Lambda$ line, where the base is contractible and the wall a point, the conservation theorem reduces to the local constancy of a signature and its content is the constancy clause on $K$ alone.
\end{enumerate}
\end{theorem}
\begin{proof}
(i) $V$ is the multiple of $w_{\Lambda}$ with $Q(V) = -3/\Lambda$ (Theorem \ref{c:odd-extensions.are-ascents-pair-signed}), and $w_{\Lambda}/\sqrt{|\Lambda|}$ has norm $\mp 1/3$, so it is $e_{V}/\sqrt 3$; the rescaled Gram's dying entry is $(-\Lambda/3)/t$ with $t = \mp\Lambda$, and the off-diagonal entries are zero for every $\Lambda$ by Proposition \ref{c:odd-extensions.line-tame-wall}(i), so the sheaf paper's limit is attained identically; the signatures are by inspection. (ii) The inverse Gram's dying entry is $-3/\Lambda$, so $\Lambda$ times it is $-3$; the sheaf paper's residue formula gives $\lim \Lambda\, G^{-1} = (1/T)\,w_{0} \otimes w_{0}$ with $T = -1/3$; and $Q(V) = -3/\Lambda$ is Fact \ref{c:odd-extensions.imported-frozen-point}. The weight statement is the sheaf paper's reading of $L^{\otimes 2} \cong N^{*}\Sigma$, and $H^{2} = Q(I) = \Lambda/3$ is \cite[Fact~2.5]{ButlerT12} (\ref{c:standard-model.imported-frozen-point-dictionary}), the form the Standard-Model paper uses for its $N = M_{P}/H$ \cite[Obs.~7.1]{ButlerT12}. (iii) The rung's metric has signature $(2,4)$ for every $\Lambda$; at $\Lambda = 0$ the five vectors of $I_{0}^{\perp}$ together with $I_{0}$ span only five dimensions, $I_{0}$ being the radical, and the sixth direction is $\varphi$ (Theorem \ref{c:odd-extensions.transport-constrains-frozen-point-only-through}(i)); for $\Lambda \neq 0$ the sixth direction is $I_{\Lambda}$; the rank count $5 + 1 = 6$ is the projective-to-conformal paper's $m + r$ (Fact \ref{c:odd-extensions.imported-minimal-completion-degenerate-algebra-rule}). For the general model, $Q(u + \tfrac{t}{2}v) = t$, the complement is spanned by $u - \tfrac{t}{2}v$ of norm $-t$, and the total is the fixed $(p+1,q+1)$; the $\Lambda$ family is the case $u = I_{0}$, $v = \varphi$ with $\ip{I_{0}}{\varphi} = 1$ and $t = -\Lambda/3$ up to the gauge. (iv) $Q(I_{\Lambda}) = \Lambda/3$ changes sign at the wall; Fact \ref{c:odd-extensions.imported-tame-walls-exceptional-algebra-residue}'s conservation theorem requires $\operatorname{sig} K$ constant across it, and its proof states that the clause is necessary because a completion allowed to switch $K$'s signature can match the two sides, which is what happens here. The script verifies (i) to (iv) on the same rational instances: the rescaled dying entries, the two exceptional signatures, $\Lambda\,(G^{-1})_{ww} = -3$ for every $\Lambda \neq 0$, the total signature $(2,4)$ on both sides with the sign of $K$ switching, the five-dimensional span at $\Lambda = 0$, and the signature $(2,4)$ of the Lorentz frame with $I_{0}$ and $\varphi$.

\end{proof}

\begin{corollary}[Two classes of completion]\label{c:odd-extensions.two-classes-completion}
\changedmark{gray}{reworded}
The sheaf paper's conservation theorem concerns completions with one exterior model: the same nondegenerate complement $K$ on both sides of the wall. For those, degeneracy is conserved and displaced to a shadow wall, and the ambient adjoins a hyperbolic plane to the whole bundle: rank two by Fact \ref{c:odd-extensions.imported-minimal-completion-degenerate-algebra-rule}, and to the whole bundle because off the wall the bundle is nondegenerate and the standard complement must exist there. The $\Lambda$ family is completed by a construction of the other class: a family of hyperplanes in one fixed nondegenerate space, whose complement is a single line that changes type across the wall, rank one. Nothing prevents a spacetime wall from being completed in the second class as well, with a complement that is not standard; what the sheaf paper proves is that the first class cannot be made two-sided. The transport principle's founding move, resolving a nilpotent by a pair, is the first class, and the frozen point's passage from anti-de Sitter to de Sitter is the second, which is not a transport (Definition \ref{c:obj.odd-extensions.single-extension}).
\end{corollary}
\begin{proof}
Theorem \ref{c:odd-extensions.exceptional-algebra-residue-completion-both-sides}(iii) and (iv) against Fact \ref{c:odd-extensions.imported-tame-walls-exceptional-algebra-residue}. The rank two of the first class rests on Fact \ref{c:odd-extensions.imported-minimal-completion-degenerate-algebra-rule}, which supplies that a minimal completion of a radical of dimension $r$ adjoins exactly $2r$ dimensions, one hyperbolic plane for a corank-one wall.

\end{proof}

\begin{remark}[what the section adds]\label{c:note.odd-extensions.what-section-adds}
That the family of intermediate algebras in $\Lambda$ is a tame wall at all is generic (Theorem \ref{c:odd-extensions.exceptional-algebra-residue-completion-both-sides}(iii)). What is specific to the cosmological constant is which objects of the dark-energy paper play the tame structure's roles: the frozen point's null direction is the radical, the Stelle-West line is the dying direction and the Stelle-West norm is the residue; the completion on both sides that the rung provides is the generic one for a family of hyperplanes and is recorded as an instance of the case the conservation theorem excludes, not as a property of $\Lambda$. None of this fixes $\Lambda$, its sign or its magnitude: the family is parametrised by $\Lambda$ and the theory describes the family.
\end{remark}

\section{What this note depends on}

\phantomsection\label{c:note.odd-extensions.setting-what-note-depends}
The note makes no prediction; no observation can refute its identifications, and the only observational input it uses is the sign of $\Lambda$, which selects which branch describes this universe (Remark \ref{c:note.odd-extensions.what-says-what-it-does-not}). What it claims is a set of identifications, each proved in the note from imported results, and this section says which imported result each identification depends on and what falls with it.
\begin{enumerate}
\item[(D1)] \emph{The frozen point.} Theorem \ref{c:odd-extensions.are-ascents-pair-signed}, Proposition \ref{c:odd-extensions.translation-bivectors-frozen-point} and all of the section ``The soldering decides'', the section ``Running the transport on the wall'', the section ``Two negative results'' and the section ``The wall in the cosmological constant'' rest on the dark-energy paper's frozen point (Fact \ref{c:odd-extensions.imported-frozen-point}): the parallel scale tractor of norm $\Lambda/3$ with the Stelle-West vector as the fifth direction of $I^{\perp}$. Without it the algebra of the section ``The frozen point splits the ascent's pair'' is a statement about an arbitrary vector of the rung's hyperbolic plane with no physical referent.
\item[(D2)] \emph{The soldered residue.} Proposition \ref{c:odd-extensions.belongs-tower} and Corollaries \ref{c:odd-extensions.two-readings-which-which} and \ref{c:odd-extensions.not-coupling-between-towers} rest on the dark-matter paper's definition of the residue as the signature of the tangent image (Fact \ref{c:odd-extensions.imported-copy-theorem-soldered-residue}). A definition admitting unsoldered directions into the residue would reopen both questions the section closes and, by the same stroke, the exactness of the dark-matter paper's superselection.
\item[(D3)] \emph{The copy theorem's inputs.} Corollary \ref{c:odd-extensions.copy-theorems-blind-spot} and Proposition \ref{c:odd-extensions.two-towers-one-curved-universe}(ii) rest on every input of the copy theorem being a spin-group object (Fact \ref{c:odd-extensions.imported-copy-theorem-soldered-residue}), and Remark \ref{c:note.odd-extensions.one-operator-sees-difference}'s conclusion that the mirror sector is observationally identical to this one rests on Fact \ref{c:dark-energy.point-infinity-has-no-physical-fluctuation}: a derived quantity depending on the bare fifth generator would need a physical coupling to the pure-gauge direction, which that Fact excludes.
\item[(D4)] \emph{The Goldstone proposition.} Proposition \ref{c:odd-extensions.boost-dilation-ratio-gauge}(iii) rests on the dark-energy paper's $(I,V)$ sector being pure gauge (Fact \ref{c:dark-energy.point-infinity-has-no-physical-fluctuation}); a dynamical fixing of the dilation gauge would make the ratio $b/a$ physical.
\item[(D5)] \emph{The residue of a conformal class.} Remark \ref{c:note.odd-extensions.holographic-coincidence}'s statement that the boundary of $\mathrm{AdS}_{4}$ has residue $7$ and lies in a tower distinct from the bulk's rests on Fact \ref{c:odd-extensions.imported-conformal-completion}, on Fact \ref{c:odd-extensions.imported-towers-clock}, and on the signature of a conformal class being well defined. A definition of the residue requiring a chosen scale would withdraw it, and with it the remark's statement of the difference between the two frameworks.
\item[(D6)] \emph{The sheaf paper's hypotheses.} the section ``The wall in the cosmological constant'' rests on the sheaf paper's tameness conditions and conservation theorem as stated (Fact \ref{c:odd-extensions.imported-tame-walls-exceptional-algebra-residue}); its one observation is about the hypothesis of that theorem and stands or falls with the theorem's statement.
\end{enumerate}
the section ``What a single generator does'' to the section ``The clock, and the paths between the oriented towers'' prove their algebra in the note, importing Facts \ref{c:odd-extensions.imported-representatives-dimension-nearest-four-types} and \ref{c:odd-extensions.imported-towers-clock} for it; their identifications with the frozen point, in Propositions \ref{c:odd-extensions.translation-bivectors-frozen-point} and \ref{c:odd-extensions.four-extensions-are-four-odd-towers}, rest on Facts \ref{c:odd-extensions.imported-conformal-completion} and \ref{c:odd-extensions.imported-frozen-point} as (D1) records, and the copy theorem's inputs in Theorem \ref{c:odd-extensions.twistor-algebra}(iii), with the superselection read in Remarks \ref{c:note.odd-extensions.what-says-what-it-does-not} and \ref{c:note.odd-extensions.no-mechanism}, on Fact \ref{c:odd-extensions.imported-copy-theorem-soldered-residue}.

\section{Position in the series}

\phantomsection\label{c:note.odd-extensions.setting-position-series}
This paper is T14a of the series, a companion to the dark-matter paper T14 \cite{ButlerT14}, downstream of its Theorem 3.4 (types by residue), Theorem 3.5 (the soldered residue), Proposition 4.1 (superselection) and Theorem 4.3 (the mirror copy), and it says where that copy theorem stops: at the real form. It uses the tower paper T3 \cite{ButlerT3} for the clock, the transport paper T1 \cite{ButlerT1} for why a single generator is not a transport, the twistor paper T9 \cite{ButlerT9} and the Standard-Model paper T12 \cite{ButlerT12} for the conformal completion and the twistor module, and the dark-energy paper T10 \cite{ButlerT10} for the frozen point, whose rank-five bundle $I^{\perp}$ is the physical occurrence of the single extension. A sibling companion note to T14 constructs the odd-dimensional twistor spaces and the odd towers' content under the reading in which all five generators are soldered; this note is upstream of nothing there and is consistent with it, the two reading the same odd algebras under different solderings, which Proposition \ref{c:odd-extensions.belongs-tower} assigns to different towers. Theorem \ref{c:odd-extensions.are-ascents-pair-signed} is downstream of T10's Theorem 3.1 and gives its parabolic-to-reductive rotation an algebraic form (Remark \ref{c:note.odd-extensions.what-adds-dark-energy-paper}); Theorem \ref{c:odd-extensions.transport-constrains-frozen-point-only-through} runs the completion theorem of the projective-to-conformal paper T2 \cite{ButlerT2} on the wall of Theorem \ref{c:odd-extensions.are-ascents-pair-signed}(iii) and finds the cosmological constant to be the one modulus the transport leaves free, which the capstone should record beside T10's parameter-of-the-laws finding; the section ``Two negative results'' adds that the frame's components are gauge by T10's Goldstone proposition and that the sign of $\Lambda$ does not select the visible tower. the section ``The wall in the cosmological constant'' runs the sheaf paper T4 \cite{ButlerT4} on the family over the $\Lambda$ line and finds it a tame wall completed in the class the conservation theorem excludes (Corollary \ref{c:odd-extensions.two-classes-completion}), which T4's successor-problem list should carry. The capstone T17 \cite{ButlerT17} should carry Corollary \ref{c:odd-extensions.copy-theorems-blind-spot} as the statement of what the copy theorem does and does not identify above the shadow, Corollary \ref{c:odd-extensions.cosmological-constant-deforms-wall} beside T10's entry, and Corollary \ref{c:odd-extensions.not-coupling-between-towers} beside the odd towers' missing coupling rule, which this note leaves as missing.

\section{Machine verification}

\phantomsection\label{c:note.odd-extensions.setting-machine-verification}
Every computational identity asserted in this note has been verified in exact rational arithmetic by the script \texttt{verify\_T14a.py} \cite{ButlerVerifyT14ext}, which imports the series' shared library \texttt{GA.py} and uses its basis-blade engine; no floating point anywhere. The script indexes generators from zero, so its $e_{0}, \ldots, e_{5}$ are the text's $e_{1}, \ldots, e_{6}$. Eighty-six checks, all passing, in nine parts.

\emph{Part A: the volume element.} Centrality and the sign $\omega^{2} = (-1)^{n(n-1)/2+q}$ on every diagonal signature of dimensions one, three, five and seven; the four signs of Corollary \ref{c:odd-extensions.four-extensions-oriented-pair}; commutation of $\omega$ with the added generator; Theorem \ref{c:odd-extensions.doubling-or-complexification}(iii) and the identities behind (i) and (ii), with the specific sign assignment of (i), over every odd signature of dimensions three and five; and the dimension of the centre, two over every odd signature of dimension at most five and one over every even signature of dimension at most six.

\emph{Part B: doubling and complexification.} For the two doubled extensions, the central idempotents $(1 \pm \omega)/2$, the two ideals of dimension sixteen, the action of $e_{5}$ as $+s\,\omega'$ and $-s\,\omega'$ on them with $s = e_{5}^{2}$, and the shadow meeting the even subalgebra exactly in its even part; for the two complexified extensions, the spanning of all thirty-two dimensions by $\Cl(p,q) + \omega\,\Cl(p,q)$ and $\omega^{2} = -1$ central.

\emph{Part C: where the mirror sits.} For all four extensions, the pairwise anticommutation and the squares of the bivectors $e_{\mu}e_{5}$, the generated subalgebra equal to the even subalgebra, the shadow of dimension sixteen meeting it in dimension eight, and the squares of $\omega e_{\mu}$ (reversed in the complexified cases, unchanged in the doubled ones); for $\Cl(2,3)$, the intersection equal to $\Cl^{0}(1,3)$.

\emph{Part D: the rung.} In $\Cl(2,4)$, the signatures $(2,3)$ and $(4,1)$ generated by $e_{i}e_{6}$ and $e_{i}e_{1}$, both generating the even subalgebra of dimension thirty-two; $\omega_{6}$ even, $\omega_{6}^{2} = -1$, commuting with every even element and anticommuting with every vector; the copies of $\Cl(1,3)$ and $\Cl(3,1)$ inside the even subalgebra.

\emph{Part E: the clock and the paths.} The residues of the four extensions, of the oriented pair and of the odd towers' representatives; the residue preserved by conformal completion; the centres of the four extensions of dimension two with the signs that separate the reversed pairs, and the centre of the rung of dimension one; the enumeration of sign-change walks from $(1,3)$ to $(3,1)$, exactly one, with $(0,4)$ a dead end.

\emph{Part F: the frozen point and the ascent's pair.} On exact rational instances of $I = a\,e_{+} + b\,e_{-}$ with $Q(I) = 3$, $-3$ and $0$: the orthogonality of the Lorentz frame and $w$ to $I$, $Q(w) = -Q(I)$; the signatures $(1,4)$, $(2,3)$ and $(1,3;1)$ of $I^{\perp}$ by exact congruence; $w = I$ and $I^{2} = 0$ in the null case; the opposite signs and anticommutation of $I$ and $w$; the generation of all sixty-four dimensions by $I^{\perp}$ with $I$ when $Q(I) \neq 0$; in $\Cl(1,3;1)$ the nilpotent generator, its anticommutation with the Lorentz frame, the sixteen-dimensional ideal and its vanishing square; and the signature $(1,3)$ of the Lorentz frame in all three cases.

\emph{Part G: the transport on the wall.} In the projective-to-conformal paper's model of the completion of $(1,3;1)$: the signature $(2,4)$ by exact congruence; the rank-five linear system for the image of the partner and its unique null solution $c = 0$; the norm $2ab$ of $aI + b\varphi$ on instances of both signs and zero; the invariance of the norm under rational boosts and the swap, the image staying orthogonal to the Lorentz frame; the centres of the two intermediate algebras; and the null pair $v, v^{*}$ in the root paper's normalisation, its identification with the canonical frame, and the norms $\pm\alpha\beta$ and orthogonality of $\alpha v \pm \beta v^{*}$.

\emph{Part H: the two negative results.} The canonical null frame in the rung's diagonal basis; $B = I \wedge \varphi = -e_{+}e_{-}$ with $B^{2} = 1$, commuting with the Lorentz bivectors and vectors and anticommuting with $e_{\pm}$; the rational rotor $5/4 + (3/4)B$ unit, fixing the Lorentz frame, scaling $I$ by $4$ and $\varphi$ by $1/4$, moving a frozen point along its hyperbola; the mirror's $I^{\perp}$ signatures $(4,1)$ and $(3,2)$ on the same instances; the two-by-two rule; the common presentation of $\Cl^{0}(1,4)$ and $\Cl^{0}(4,1)$; and the squares and factorisations of the bare fifth generator.

\emph{Part I: the wall in the cosmological constant.} In four dilation gauges and at seven rational values of $\Lambda$ including zero: the Gram $\operatorname{diag}(1,-1,-1,-1,-\Lambda/3)$ of $I_{\Lambda}^{\perp}$ and its corank at $\Lambda = 0$ only; the gauge independence of the dying eigenvalue; the exact crossing slope $-1/3$ in the parametrisation by $\Lambda$; the signatures $(2,3)$ and $(1,4)$ on the two sides; the rescaled dying entries $\pm 1/3$ and the exceptional signatures; $\Lambda\,(G^{-1})_{ww} = -3$ in every gauge with $w_{0} = a\,I_{0}$; the total signature $(2,4)$ on both sides with the complement's sign switching; the five-dimensional span at the wall and the signature of the Lorentz frame with $I_{0}$ and $\varphi$; and the general hyperplane model of Theorem \ref{c:odd-extensions.exceptional-algebra-residue-completion-both-sides}(iii) for five signatures $(p,q)$ and five values of $t$.

What the script does not check: the isomorphism types $\Mat(2,\HH)$ and $\Mat(4,\R)$ of the two oriented shadows (Fact \ref{c:odd-extensions.imported-representatives-dimension-nearest-four-types}, imported; the note uses that they differ and, in Corollary \ref{c:odd-extensions.copy-theorems-blind-spot}, which type belongs to which tower), the frozen-point geometry (Facts \ref{c:odd-extensions.imported-frozen-point} and \ref{c:odd-extensions.imported-frozen-point-degenerate-algebras}, imported; the script checks the algebra of a vector of each norm sign in the rung, not that the parallel tractor is that vector), the algebra-map properties of Theorem \ref{c:odd-extensions.doubling-or-complexification}(i) and (ii) (argued in the text; the script checks the identities they rest on and the four extensions in detail), the identification of the MacDowell-Mansouri translation part with the bivectors $e_{\mu}e_{5}$ beyond the signs of their squares, the reversal of the tractor norm under $g \mapsto -g$ (Proposition \ref{c:odd-extensions.two-towers-one-curved-universe}(i), argued in the text; Part H checks its consequences from the reversed signature declared as input), and the transitivity of the boosts on each level set of the norm (Theorem \ref{c:odd-extensions.transport-constrains-frozen-point-only-through}(iii), argued in the text; Part G checks invariance under particular boosts).

\section{Results used from elsewhere}
The proofs above use the following results of the series, stated briefly; each is proved in the paper named.

\begin{fact}[Minimal nondegenerate completion]\label{c:conformal.minimal-nondegenerate-completion}
A real quadratic space with radical of dimension $r$ and nondegenerate part of signature $(p,q)$ has minimal nondegenerate completions of exactly $r$ extra dimensions, all of signature $(p+r,q+r)$ and unique up to an isometry fixing the space, so $\Cl(p,q,r)$ embeds in $\Cl(p+r,q+r)$, minimally among completion-induced embeddings.
\emph{Source:} T2, Thm 2.3 (\href{https://doi.org/10.5281/zenodo.22714886}{doi:10.5281/zenodo.22714886}).
\end{fact}

\begin{fact}[Imported: the forced completion and its fibre quadric]\label{c:correspondence.imported-forced-completion-fibre-quadric}
Division by a type-N zero divisor in $\Cl(1,3)$ has no element-level resolution, its unique minimal resolving completion is $\Cl(1,3) \hookrightarrow \Cl(2,4)$, and the resolving assignments lie on the null quadric of $\R^{2,4}$.
\emph{Source:} T9, Fact 2.1 (\href{https://doi.org/10.5281/zenodo.22722290}{doi:10.5281/zenodo.22722290}).
\end{fact}

\begin{fact}[The spinor identification]\label{c:correspondence.spinor-identification}
There is a canonical $\Spin$-equivariant isomorphism from $V_\C$ to the bivectors of the conformal spinor space carrying the quadratic form to the wedge square, so null vectors are exactly the simple bivectors, that is, the lines of the projective spinor space.
\emph{Source:} T9, Thm 3.4 (\href{https://doi.org/10.5281/zenodo.22722290}{doi:10.5281/zenodo.22722290}).
\end{fact}

\begin{fact}[The frozen point makes the Stelle-West vector]\label{c:dark-energy.frozen-point-makes-stelle-west-vector}
On the Einstein region with $\Lambda\neq 0$ the frozen point reduces the ambient structure to de Sitter or anti-de Sitter type and, with the canonical tractor, supplies the Stelle-West vector $V = X - (3/\Lambda)I$, so the MacDowell-Mansouri action is $Q(I)^{-1}\int\epsilon(I,V,F,F)$, polynomial and free of inverse metrics.
\emph{Source:} T10, Thm 3.1 (\href{https://doi.org/10.5281/zenodo.22722496}{doi:10.5281/zenodo.22722496}).
\end{fact}

\begin{fact}[Imported: the frozen point and its geometry]\label{c:dark-energy.imported-frozen-point-geometry}
On the invertible region of a gravitational field on flat spacetime the Einstein condition makes the scale tractor parallel for the tractor gauge field, with $Q(I) = \Lambda/3$, stabiliser $SO(1,4)$ or $SO(2,3)$, and $P_{ab} = (\Lambda/6)g_{ab}$.
\emph{Source:} T10, Fact 2.1 (\href{https://doi.org/10.5281/zenodo.22722496}{doi:10.5281/zenodo.22722496}).
\end{fact}

\begin{fact}[The point at infinity has no physical fluctuation: the $(I,V)$-sector is pure gauge]\label{c:dark-energy.point-infinity-has-no-physical-fluctuation}
\changedmark{gray}{reworded}
Within the three-argument variational principle, the constrained fluctuations of $(I,V)$ and the broken generators of $SO(2,4)$ both number nine and correspond bijectively, so the $(I,V)$-sector is pure gauge and the graviton has no $(\delta I,\delta V)$ polarisation.
\emph{Source:} T10, Prop 5.4 (\href{https://doi.org/10.5281/zenodo.22722496}{doi:10.5281/zenodo.22722496}).
\end{fact}

\begin{fact}[Imported: the eight towers and the mod-8 clock]\label{c:dark-matter.imported-eight-towers-mod-8-clock}
The $(1,1)$ ascent preserves $p - q$ with $\Cl(p+k,q+k) \cong \Mat(2^{k}, \Cl(p,q))$, so the eight residues of $p - q$ modulo eight label eight parallel towers.
\emph{Source:} T14, Fact 2.5 (\href{https://doi.org/10.5281/zenodo.22722953}{doi:10.5281/zenodo.22722953}).
\end{fact}

\begin{fact}[The mirror tower is a copy of this one]\label{c:dark-matter.mirror-tower-copy-one}
Every input of the Standard-Model chain is invariant under signature reversal, so the mirror tower carries a copy of the derived Standard Model and, with a common ceiling $M$ and under the single frozen point and indifference (which discharge the democracy premise for the mirror), $\Omega_{\mathrm{mirror}} = \Omega_{b}$ to half a percent; T14c finds that indifference with this reach would also make the mirror as hot as ours, so the cold mirror is an exception to it.
\emph{Source:} T14, Thm 4.3 (\href{https://doi.org/10.5281/zenodo.22722953}{doi:10.5281/zenodo.22722953}).
\end{fact}

\begin{fact}[Projection: every tower's content dilutes as this one's]\label{c:dark-matter.projection-every-towers-content-dilutes-ones}
Under the sourcing rule read from P1, with any extra directions co-expanding or static and a missing direction filled by comoving sheets, every tower's projected matter dilutes as $a^{-3}$ and its radiation as $a^{-4}$, as in this tower.
\emph{Source:} T14, Lem 6.1 (\href{https://doi.org/10.5281/zenodo.22722953}{doi:10.5281/zenodo.22722953}).
\end{fact}

\begin{fact}[The residue is a soldered invariant, not an algebraic one]\label{c:dark-matter.residue-soldered-invariant-not-algebraic-one}
The residue is preserved by every isomorphism carrying vectors to vectors, but not by abstract algebra isomorphisms, under which residues $0$ and $2$, $3$ and $7$, and $4$ and $6$ coincide at equal dimension.
\emph{Source:} T14, Thm 3.5 (\href{https://doi.org/10.5281/zenodo.22722953}{doi:10.5281/zenodo.22722953}).
\end{fact}

\begin{fact}[Superselected, hence dark and exactly stable]\label{c:dark-matter.superselected-hence-dark-exactly-stable}
Under P1, the sector decomposition and the wall-sequence theorem's hypotheses (compact base, good couplings), the tower charge is exactly superselected, so dark-tower matter is exactly stable; that it only gravitates, at first order, is the series' acoustic no-go's, within the s-wave current algebra and conditional on the Indifference premise and premise (Q).
\emph{Source:} T14, Prop 4.1 (\href{https://doi.org/10.5281/zenodo.22722953}{doi:10.5281/zenodo.22722953}).
\end{fact}

\begin{fact}[Towers are residues, and the single-clock filter]\label{c:dark-matter.towers-are-residues-single-clock-filter}
The residue $p - q \pmod 8$ is invariant along each tower, and under the nearest-dimension premise P3, with P2 reading the time-sign label as indexical, exactly the residues $0$ and $4$ lack a single time direction, leaving six towers that admit matter (the two-time exclusion imports the standard instability).
\emph{Source:} T14, Prop 3.1 (\href{https://doi.org/10.5281/zenodo.22722953}{doi:10.5281/zenodo.22722953}).
\end{fact}

\begin{fact}[Type by residue, and what each tower can carry intrinsically]\label{c:dark-matter.type-residue-what-each-tower-can}
A tower's spinor type (real, complex or quaternionic, with one or two modules) is a function of its residue alone, and it fixes the intrinsic symmetry group, $\Z_{2}$, $U(1)$ or $Sp(1)$, that each tower's ground algebra carries.
\emph{Source:} T14, Thm 3.4 (\href{https://doi.org/10.5281/zenodo.22722953}{doi:10.5281/zenodo.22722953}).
\end{fact}

\begin{fact}[The sign of $\omega^2$]\label{c:primer.sign}
The pseudoscalar of $\Cl(p,q)$ satisfies $\omega^{2} = (-1)^{(p-q)(p-q-1)/2}$, so $\omega^{2} = +1$ exactly when $p - q \equiv 0$ or $1 \pmod 4$.
\emph{Source:} T0, Lem 9.5 (\href{https://doi.org/10.5281/zenodo.22714501}{doi:10.5281/zenodo.22714501}).
\end{fact}

\begin{fact}[Conservation of degeneracy, absolute form]\label{c:sheaf.conservation-degeneracy-absolute-form}
Assuming a tame signature-change wall on a connected manifold whose two sides are nonempty, every completion that is standard, with one exterior model, outside a neighbourhood of the wall degenerates somewhere.
\emph{Source:} T4, Thm 5.1 (\href{https://doi.org/10.5281/zenodo.22721669}{doi:10.5281/zenodo.22721669}).
\end{fact}

\begin{fact}[The degeneration sheaf]\label{c:sheaf.degeneration-sheaf}
For a corank-one degeneration satisfying the tameness hypotheses, the Clifford fibre at each wall point carries a canonical square-zero ideal of dimension $2^{n-1}$ generated by the radical, with quotient the Clifford algebra of the induced nondegenerate form, and these ideals assemble into sheaves supported exactly on the wall.
\emph{Source:} T4, Thm 3.1 (\href{https://doi.org/10.5281/zenodo.22721669}{doi:10.5281/zenodo.22721669}).
\end{fact}

\begin{fact}[The exceptional algebra: deformation to the normal cone]\label{c:sheaf.exceptional-algebra-deformation-normal-cone}
Under tameness, rescaling the dying direction by $1/\sqrt{t}$ along a transversal gives the nondegenerate limit form $G_0|_{V_1} \perp \dot G(k,k)$, of signature $(p,q)$ on the physical side, whose Clifford algebra is canonically independent of the choices.
\emph{Source:} T4, Prop 6.1 (\href{https://doi.org/10.5281/zenodo.22721669}{doi:10.5281/zenodo.22721669}).
\end{fact}

\begin{fact}[The invariant form of (T3)]\label{c:sheaf.invariant-form-t3}
Given a smooth wall with a smooth corank-one radical line, the differential of $G(k,k)$ along the wall defines a canonical bundle map $\delta : \shL^{\otimes 2} \to N^*\Sig$, and transversality is exactly the demand that it be an isomorphism.
\emph{Source:} T4, Lem 2.2 (\href{https://doi.org/10.5281/zenodo.22721669}{doi:10.5281/zenodo.22721669}).
\end{fact}

\begin{fact}[The metric's residue, and the square-root twist]\label{c:sheaf.metrics-residue-square-root-twist}
Assuming the wall is tame, the metric has a canonical simple pole whose residue is the inverse of the transversality datum, giving $N^*\Sig \cong \shL^{\otimes 2}$ and making the completion's partner line a square root of the normal bundle.
\emph{Source:} T4, Thm 6.2 (\href{https://doi.org/10.5281/zenodo.22721669}{doi:10.5281/zenodo.22721669}).
\end{fact}

\begin{fact}[One-sided completion]\label{c:sheaf.one-sided-completion}
Under tameness, with signature $(p,q)$, $p \geq 1$, on the physical side and a trivial radical line bundle, coupling an adjoined hyperbolic plane to the dying direction gives an ambient form of constant signature $(p{+}1, q{+}1)$ on the whole physical side, wall included.
\emph{Source:} T4, Thm 4.4 (\href{https://doi.org/10.5281/zenodo.22721669}{doi:10.5281/zenodo.22721669}).
\end{fact}

\begin{fact}[Resolution, ideal-theoretically]\label{c:sheaf.resolution-ideal-theoretically}
In any nondegenerate quadratic space containing a wall fibre with its form, the radical line generates the unit ideal of the Clifford algebra.
\emph{Source:} T4, Prop 3.4 (\href{https://doi.org/10.5281/zenodo.22721669}{doi:10.5281/zenodo.22721669}).
\end{fact}

\begin{fact}[Imported: the Fock representation of a Clifford algebra, and its Bogoliubov group]\label{c:standard-model.imported-fock-representation-clifford-algebra-bogoliubov}
The Fock space $\Lambda^{\bullet}E$ of $n$ fermionic modes is the irreducible spinor module of the complex Clifford algebra on $2n$ generators, with $\mathrm{Pin}(2n)$ implementing the Bogoliubov group and $U(n)$ the number-preserving part.
\emph{Source:} W. Fulton and J. Harris, Representation Theory: A First Course, Graduate Texts in Mathematics 129, Springer (1991)., Lecture~20, \S 20.1; H. B. Lawson and M.-L. Michelsohn, Spin Geometry, Princeton University Press (1989)., Ch.~I, \S\S 2, 5.
\end{fact}

\begin{fact}[Imported: the frozen point and the dictionary]\label{c:standard-model.imported-frozen-point-dictionary}
Mass is conformal weight read against the frozen point, $m^{2} = \Delta(3-\Delta)H^{2}$, so heavy fields have complex weight and the Planck mass's imaginary weight is $N \approx M_{P}/H$, about $10^{61}$.
\emph{Source:} T12, Fact 2.5 (\href{https://doi.org/10.5281/zenodo.22722669}{doi:10.5281/zenodo.22722669}).
\end{fact}

\begin{fact}[Exponentials of bivectors with scalar square]\label{c:std-algebra.bivector-exponential}
If $A^{2} = \lambda \in \R$ then $e^{A}$ is $\cos\theta + \frac{\sin\theta}{\theta}A$, $\cosh\theta + \frac{\sinh\theta}{\theta}A$ or $1 + A$, a rotor with $e^{A}\widetilde{e^{A}} = 1$.
\emph{Source:} the algebra core of the Transport Lab, proved there with an exact check; standard, see C. Doran and A. Lasenby, Geometric Algebra for Physicists, Cambridge University Press (2003), Ch. 4 (foundations of geometric algebra: rotors, versors).
\end{fact}

\begin{fact}[Orthogonal vectors anticommute]\label{c:std-algebra.orthogonal-anticommute}
$uv = -vu$ iff $B(u,v) = 0$; in an orthogonal basis $e_I e_J = \pm \prod_{I \cap J} Q(e_i)\; e_{I \triangle J}$ with the sorting sign.
\emph{Source:} the algebra core of the Transport Lab, proved there with an exact check; standard, see C. Doran and A. Lasenby, Geometric Algebra for Physicists, Cambridge University Press (2003), Ch. 4 (foundations of geometric algebra).
\end{fact}

\begin{fact}[The Pauli algebra: Cl(3,0) and 2 x 2 complex matrices]\label{c:std-algebra.pauli-algebra}
$\Cl(3,0) \cong \Mat_2(\C)$ with $e_k \mapsto \sigma_k$, and the pseudoscalar $e_1e_2e_3$ maps to $\mathrm{i}$.
\emph{Source:} the algebra core of the Transport Lab, proved there with an exact check; standard, see C. Doran and A. Lasenby, Geometric Algebra for Physicists, Cambridge University Press (2003), Ch. 2 (geometric algebra in two and three dimensions).
\end{fact}

\begin{fact}[How the pseudoscalar commutes]\label{c:std-algebra.pseudoscalar-commutation}
$I X_k = (-1)^{k(n-1)} X_k I$; $I\tilde{I} = \prod_i Q(e_i)$, and $I^{2} = 0$ in a degenerate algebra.
\emph{Source:} the algebra core of the Transport Lab, proved there with an exact check; standard, see C. Doran and A. Lasenby, Geometric Algebra for Physicists, Cambridge University Press (2003), Ch. 4 (foundations of geometric algebra).
\end{fact}

\begin{fact}[The spacetime algebra as a matrix algebra, and the spacetime split]\label{c:std-algebra.sta-matrix}
$\Cl(1,3) \cong \Mat_2(\HH)$, the Dirac matrices represent $\C \otimes \Cl(1,3) \cong \Mat_4(\C)$, and $\sigma_k = \gamma_k\gamma_0$ give $\Cl^{+}(1,3) \cong \Cl(3,0)$ with $\sigma_1\sigma_2\sigma_3 = I$.
\emph{Source:} the algebra core of the Transport Lab, proved there with an exact check; standard, see C. Doran and A. Lasenby, Geometric Algebra for Physicists, Cambridge University Press (2003), Ch. 5 (relativity and spacetime).
\end{fact}

\begin{fact}[The even subalgebra and the Z2 grading]\label{c:std-algebra.z2-grading}
$\Cl = \Cl^{+} \oplus \Cl^{-}$ is a $\mathbb{Z}_2$-grading, $\dim \Cl^{+} = 2^{n-1}$, and $\Cl^{j}\Cl^{k}$ has grades $|j-k|, |j-k|+2, \dots, j+k$.
\emph{Source:} the algebra core of the Transport Lab, proved there with an exact check; standard, see P. Lounesto, Clifford Algebras and Spinors, 2nd ed., London Mathematical Society Lecture Note Series 286, Cambridge University Press (2001), Ch. 14 (definitions of the Clifford algebra).
\end{fact}

\begin{fact}[Commuting copies]\label{c:tower.commuting-copies}
Each floor $\Cl(p{+}k, q{+}k)$ of the tower is canonically a twisted copy of $\Cl(p,q)$ tensored with $\Mat_{2^k}(\R)$, that is $\Mat_{2^k}(\Cl(p,q))$.
\emph{Source:} T3, Thm 3.2 (\href{https://doi.org/10.5281/zenodo.22721556}{doi:10.5281/zenodo.22721556}).
\end{fact}

\begin{fact}[Rigidity; the clock is frozen]\label{c:tower.rigidity-clock-frozen}
Every floor of the tower shares the ground floor's charge, division-ring type, centre, simplicity class, Morita class and obstruction taxonomy, so the physical tower $\Cl(1{+}k, 3{+}k) \cong \Mat_{2^{k+1}}(\HH)$ is frozen at clock position $6$.
\emph{Source:} T3, Thm 3.6 (\href{https://doi.org/10.5281/zenodo.22721556}{doi:10.5281/zenodo.22721556}).
\end{fact}

\begin{fact}[Minimality]\label{c:transport.minimality}
Every null-lift completion adjoins directions of both signs, so the minimal completion is the hyperbolic-plane ascent $\Cl(p,q) \to \Cl(p{+}1,q{+}1)$, unique up to isometry, with increment $(+1,+1)$.
\emph{Source:} T1, Thm 4.13 (\href{https://doi.org/10.5281/zenodo.21743863}{doi:10.5281/zenodo.21743863}).
\end{fact}

\begin{fact}[One generator is minimal; the minimal completions classified]\label{c:transport.one-generator-minimal-minimal-completions-classified}
For $p - q \equiv 1 \pmod 4$, the minimal simple Clifford completions of $\Cl(p,q)$ are exactly the two one-generator ascents $\Cl(p{+}1,q)$ and $\Cl(p,q{+}1)$, isomorphic as algebras, while a null generator only thickens the algebra.
\emph{Source:} T1, Thm 5.7 (\href{https://doi.org/10.5281/zenodo.21743863}{doi:10.5281/zenodo.21743863}).
\end{fact}

\begin{refutation}[The two one-generator ascents differ as algebras]\label{c:refuted.the-two-one-generator-ascents-differ-as-algebras}
Refuted: the two one-generator ascents are always isomorphic as algebras, the sign of the adjoined generator being a geometric bit the algebra forgets.
\end{refutation}

\begin{prediction}[Direct detection stays null]\label{c:pred.dark-matter-direct-detection-null}
Conditionally on the single-frozen-point premise, and on the series' acoustic no-go (s-wave current algebra, Indifference premise, premise (Q)), direct-detection searches should stay null permanently, because dark matter couples to ordinary matter through gravity alone at first order.
\emph{Kill condition:} Any confirmed dark-matter detection above the gravitational floor.
\end{prediction}

\begin{prediction}[No indirect dark-sector signals]\label{c:pred.no-indirect-dark-signals}
Conditionally on the single-frozen-point premise, and on the series' acoustic no-go (s-wave current algebra, Indifference premise, premise (Q)), indirect-detection surveys should see no dark-sector decay products, and no annihilation products above the gravitationally suppressed rate of a charge-asymmetric relic.
\emph{Kill condition:} Any confirmed indirect signal.
\end{prediction}

\section*{Changes since the deposited version}
\addcontentsline{toc}{section}{Changes since the deposited version}
Since the deposited version 1.0: 2 results corrected, 3 changed, 17 reworded (wording, notation or citations only) and 37 cited results added. The comparison is with the source of version 1.0 on Zenodo, \href{https://doi.org/10.5281/zenodo.23232104}{doi:10.5281/zenodo.23232104}. A result counts as revised when the words or mathematics of its statement or proof differ from that source; a citation added in a parenthesis does not count. A result the paper cites from elsewhere is compared with the deposit of the paper that proves it. Each revision is classed as corrected, changed or reworded; the classes were assigned by a review of each difference and the store's record of why it was made, and a revision not yet reviewed is counted as changed.
\subsection*{Corrected}
The deposited version was wrong or claimed too much here.
\begin{itemize}\setlength{\itemsep}{2pt}
\item Assembled~\ref{c:odd-extensions.imported-representatives-dimension-nearest-four-types} (\emph{Imported: the representatives at the dimension nearest four, and the types}), Fact 1.2 in the deposit: statement and proof. Corrected: the other six residues admit matter; the source does not say they are populated. \emph{Store notes.} \emph{2026-10-21:} The restatement said the other six residues 'are populated'; the originals (T14 Prop 3.1, Thm 3.4) say they admit matter, and which the descent populated is a further question (the odd towers' population is held as conjectured in dark-matter.which-towers-are-populated). Corrected to 'admit matter' at every level and in the proof; the paper's wording is kept [...]
\item Theorem~\ref{c:odd-extensions.doubling-or-complexification} (\emph{Doubling or complexification}), Thm 2.3 in the deposit: statement and proof. Corrected sign: on the summand of (1 plus or minus omega)/2, e\_n acts as plus or minus s omega', with s the square of e\_n. \emph{Store notes.} \emph{2026-10-26:} Grounding pass 2026-10-26: cites std-algebra.z2-grading, std-algebra.orthogonal-anticommute where the proof uses them.
\end{itemize}
\subsection*{Changed}
The content changed by a later decision, not because the deposited version was wrong.
\begin{itemize}\setlength{\itemsep}{2pt}
\item Assembled~\ref{c:odd-extensions.imported-frozen-point} (\emph{Imported: the frozen point}), Fact 1.5 in the deposit: statement and proof. Gauge Theory Gravity restatement, and the translation part renormalised to the store convention: -e wedge V over l squared. \emph{Store notes.} \emph{2026-10-23:} Restated in Gauge Theory Gravity: the invertible region is that of a gravitational field on flat spacetime, the tractor metric the square of Cl(2,4), and the soldering form the physical frame hb\^{}\{-1\}(a), so that the reduced connection is Omega(a) - hb\^{}\{-1\}(a) \^{} V/l\^{}2 in the store normalisation; content unchanged. \emph{2026-10-23:} Squares of vectors written Q(v) in place of the inner product of a vector with itself; no change of content.
\item Corollary~\ref{c:odd-extensions.four-extensions-oriented-pair} (\emph{The four extensions of the oriented pair}), Cor 2.4 in the deposit: statement. New remark: the labels follow a uniform convention; in T14's causal convention the mirror's extensions are Cl(4,1) and Cl(3,2). \emph{Store notes.} \emph{2026-10-26:} Grounding pass 2026-10-26: cites std-algebra.sta-matrix where the proof uses them.
\item Corollary~\ref{c:odd-extensions.copy-theorems-blind-spot} (\emph{The copy theorem's blind spot}), Cor 4.2 in the deposit: statement and proof. Twistor module renamed conformal spinor module; new clause names the complex unit as the conformal pseudoscalar I6, and the proof now spells out its sources. \emph{Store notes.} \emph{2026-10-21:} Converted the twistor module to the conformal spinor module $\CS$ and named the complex unit of $\Cl^0(2,4)$ as $I_6$ (TWGA E1.2). Kept the complexified ambient, since the point is that the vector space has no signature over $\C$; the proof is a citation and is unchanged.
\end{itemize}
\subsection*{Reworded}
Wording, notation, terminology or citations only; the mathematics and the claim are as deposited.
\begin{itemize}\setlength{\itemsep}{2pt}
\item Assembled~\ref{c:odd-extensions.imported-copy-theorem-soldered-residue} (\emph{Imported: the copy theorem, and the soldered residue}), Fact 1.3 in the deposit: statement and proof. Terminology ('conformal spinor module', reversion form) and citations; the fact's text now repeated as its proof.
\item Assembled~\ref{c:odd-extensions.imported-conformal-completion} (\emph{Imported: the conformal completion}), Fact 1.4 in the deposit: statement and proof. Terminology ('conformal spinor space' with complex unit I\_6) and citations only.
\item Theorem~\ref{c:odd-extensions.mirror-inside-complexified-extension-absent-doubled} (\emph{The mirror inside the complexified extension, and absent from the doubled one}), Thm 3.1 in the deposit: proof. Cross-references and citations of standard algebra facts only.
\item Proposition~\ref{c:odd-extensions.translation-bivectors-frozen-point} (\emph{The translation bivectors of the frozen point}), Prop 3.2 in the deposit: proof. Notation only: Q(V) for the square of a vector; cross-references updated.
\item Theorem~\ref{c:odd-extensions.twistor-algebra} (\emph{The even conformal algebra is $\Cl(2,3)$}), Thm 4.1 in the deposit: statement and proof. Terminology ('conformal spinor space'), with the volume element glossed as I\_6 and citations added; mathematics unchanged.
\item Proposition~\ref{c:odd-extensions.four-extensions-are-four-odd-towers} (\emph{The four extensions are the four odd towers}), Prop 5.1 in the deposit: proof. Cross-references only: citations updated and the sources of each step named in the proof.
\item Assembled~\ref{c:odd-extensions.imported-frozen-point-degenerate-algebras} (\emph{Imported: the frozen point at $\Lambda = 0$, and the degenerate algebras}), Fact 6.1 in the deposit: statement and proof. Notation only: Q(v) for the square of a vector; provenance restated with citations.
\item Theorem~\ref{c:odd-extensions.are-ascents-pair-signed} (\emph{$I$ and $V$ are the ascent's pair, signed by $\Lambda$}), Thm 6.2 in the deposit: statement and proof. Notation only: Q(v) for the square of a vector; cross-references updated.
\item Corollary~\ref{c:odd-extensions.not-coupling-between-towers} (\emph{$\Lambda$ is not a coupling between towers}), Cor 7.3 in the deposit: proof. Proof now cites the theorem supplying Lambda as a datum of the ascent; the argument is unchanged.
\item Remark~\ref{c:note.odd-extensions.dictionary-ascents-null-pair} (\emph{the dictionary with the ascent's null pair}), Rem 8.4 in the deposit: statement. Notation only: Q(v) for the square of a vector; cross-references updated.
\item Proposition~\ref{c:odd-extensions.boost-dilation-ratio-gauge} (\emph{The boost is the dilation; the ratio $b/a$ is gauge}), Prop 9.2 in the deposit: proof. Citations added for standard algebra facts (orthogonal vectors anticommute, bivector exponential); proof unchanged.
\item Proposition~\ref{c:odd-extensions.two-towers-one-curved-universe} (\emph{The two towers on one gravitational field}), Prop 9.3 in the deposit: statement and proof. Gauge Theory Gravity phrasing: the one gravitational field on flat spacetime replaces the one manifold; Q(v) notation; clauses unchanged.
\item Proposition~\ref{c:odd-extensions.line-tame-wall} (\emph{The $\Lambda$ line is a tame wall}), Prop 10.2 in the deposit: statement and proof. Notation only: Q(v) for the square of a vector; proof cites the relation 2ab = Lambda/3 and the tameness conditions explicitly.
\item Theorem~\ref{c:odd-extensions.exceptional-algebra-residue-completion-both-sides} (\emph{Exceptional algebra, residue, and the completion on both sides}), Thm 10.3 in the deposit: statement and proof. Notation only: Q(v) for the square of a vector; citations updated.
\item Corollary~\ref{c:odd-extensions.two-classes-completion} (\emph{Two classes of completion}), Cor 10.4 in the deposit: proof. Proof now names the Fact supplying the rank two of the first class; content unchanged.
\item Fact~\ref{c:dark-energy.point-infinity-has-no-physical-fluctuation} (\emph{The point at infinity has no physical fluctuation: the $(I,V)$-sector is pure gauge}), Fact 9.1 in the deposit: statement and proof. Notation Q(v) for squares and Gauge Theory Gravity wording for the tractor connection; the counting is unchanged.
\end{itemize}
\noindent In addition, 37 results from elsewhere in the series or the algebra core are stated here that the deposited paper cited or used without stating; 1 result gained citations only. These are not marked.
\begin{thebibliography}{99}
\bibitem{ButlerT14} L.~Butler,
\emph{Dark Matter and the Arrow of Time: The Five Dark Towers},
Transport Papers T14 (2026).
\bibitem{ButlerT3} L.~Butler,
\emph{The Tower and the CAR Limit: Iterated Transport, the Frozen Clock, and Second Quantisation},
Transport Papers T3 (2026), doi:\href{https://doi.org/10.5281/zenodo.22721556}{10.5281/zenodo.22721556}.
\bibitem{ButlerT9} L.~Butler,
\emph{The Twistor Correspondence: Ascended Spinors, Incidence, and the Palatial Bridge},
Transport Papers T9 (2026).
\bibitem{ButlerT12} L.~Butler,
\emph{The Standard Model from the Minimal Massive Composite},
Transport Papers T12 (2026), doi:\href{https://doi.org/10.5281/zenodo.22722669}{10.5281/zenodo.22722669}.
\bibitem{ButlerT10} L.~Butler,
\emph{Dark Energy and the Action of Gravity: MacDowell-Mansouri, the Frozen Point, and One Way to Grip},
Transport Papers T10 (2026), doi:\href{https://doi.org/10.5281/zenodo.22722496}{10.5281/zenodo.22722496}.
\bibitem{MacDowellMansouri1977} S.~W.~MacDowell and F.~Mansouri,
Unified geometric theory of gravity and supergravity,
\emph{Phys.\ Rev.\ Lett.} \textbf{38} (1977), 739--742.
\bibitem{StelleWest1980} K.~S.~Stelle and P.~C.~West,
Spontaneously broken de Sitter symmetry and the gravitational holonomy group,
\emph{Phys.\ Rev.\ D} \textbf{21} (1980), 1466--1488.
\bibitem{ButlerT6} L.~Butler,
\emph{The Geometry of the Black-Hole Interior: Where Division Fails, the Weighted Limit, and the Kerr Tests},
Transport Papers T6 (2026), doi:\href{https://doi.org/10.5281/zenodo.22721990}{10.5281/zenodo.22721990}.
\bibitem{ButlerT1} L.~Butler,
\emph{The Transport Principle: Total Division in Clifford Algebras via Minimal Completion},
Transport Papers T1 (2026).
\bibitem{Planck2018} N.~Aghanim et al. (Planck Collaboration),
\emph{Planck 2018 results. VI. Cosmological parameters},
Astron. Astrophys. \textbf{641} (2020), A6.
\bibitem{ButlerT4} L.~Butler,
\emph{Clifford Sheaves over Degenerate Metrics: Nilpotent Ideals, One-Sided Completion, and the Exceptional Algebra},
Transport Papers T4 (2026), doi:\href{https://doi.org/10.5281/zenodo.22721669}{10.5281/zenodo.22721669}.
\bibitem{ButlerT2} L.~Butler,
\emph{Projective Completes to Conformal: The Minimal Nondegenerate Completion of a Degenerate Clifford Algebra, and Where PGA's Missing Division Goes},
Transport Papers T2 (2026), doi:\href{https://doi.org/10.5281/zenodo.22714886}{10.5281/zenodo.22714886}.
\bibitem{ButlerT17} L.~Butler,
\emph{The Transport Programme: Method, Status, and the Falsifiable Ledger},
Transport Papers T17 (2026).
\bibitem{ButlerVerifyT14ext} L.~Butler,
\emph{verify\_T14a.py: machine verification for the companion note ``One Dimension at a Time: The Odd Extensions of the Oriented Towers''},
Python script, deposited with this note alongside the shared library GA.py (2026).
\end{thebibliography}
\end{document}